\documentclass[11pt]{article}

\usepackage[margin=1.05in]{geometry}
\usepackage[utf8]{inputenc}
\usepackage[T1]{fontenc}
\usepackage{amsmath,amssymb,amsthm,mathtools}
\usepackage{booktabs,array,longtable,tabularx}
\usepackage{graphicx}
\usepackage{enumitem}
\usepackage{xcolor}
\usepackage{microtype}
\usepackage{csquotes}
\usepackage{pdflscape}
\usepackage{placeins}

\usepackage[
  authordate,
  backend=biber,
  doi=true,
  url=true,
  isbn=false
]{biblatex-chicago}


\usepackage{hyperref}

\newcolumntype{L}[1]{>{\raggedright\arraybackslash}p{#1}}

\hypersetup{
  colorlinks=true,
  linkcolor=blue!50!black,
  citecolor=blue!50!black,
  urlcolor=blue!50!black
}

\graphicspath{{./}}

\newtheorem{definition}{Definition}
\newtheorem{proposition}{Proposition}
\newtheorem{remark}{Remark}

\newcommand{\Votes}{V}
\newcommand{\Seats}{S}
\newcommand{\Quota}{q}
\newcommand{\Diff}{d}
\newcommand{\Parties}{P}
\newcommand{\Coalition}{C}
\newcommand{\ChamberSeats}{513}
\newcommand{\MajoritySeats}{257}

% Empirical prose uses literal values and paragraph-local provenance.
% row_at_generation counts 1-based CSV data rows, excluding the header.
% Semantic keys identify observations; row numbers are navigation aids only.

\title{Seat Majorities without Vote Majorities:\\[0.3em]Coalition Inversions in Brazil's Chamber of Deputies}
\author{
Marcelo V. Maciel\\
Department of Political Science, University of California, Irvine\\
Irvine, CA, United States\\
Corresponding author: mvmaciel@uci.edu\\
ORCID: 0000-0001-7666-8494 }
\date{}


\addbibresource{refs2.bib}


\begin{document}
\maketitle

\begin{abstract}
  Proportional representation is usually evaluated party by party, by comparing
  each party's share of seats with its share of votes. Yet fragmented
  legislatures are governed through combinations of parties, so small
  party-level deviations can become consequential when parties are combined.
  Brazil provides a useful setting because the Chamber of Deputies is assembled
  from twenty-seven separate district elections. Several institutional steps
  separate national votes from final party seat totals, and no national
  compensatory allocation reconciles them. I examine which politically
  meaningful sets of parties hold a legislative majority with less than half of
  all valid federal-deputy votes nationwide, and how within-district allocation
  and the distribution of support contribute to these coalition
  inversions. Using the 2014, 2018, and 2022 elections, I compare distinct
  election-year cabinet party sets with potential coalitions restricted by
  ideological proximity. Cabinet inversions appear following the 2014 and 2022
  elections, but not the 2018 election; inversions among minimal connected
  winning coalitions appear in all three. In 2022, a seven-party bloc of
  ideologically adjacent parties combined 45.26 percent of the vote with 50.29
  percent of the seats. Coalition formation turns otherwise continuous vote-seat
  deviations into a discontinuous political consequence when their accumulation
  carries a party set across the legislative majority threshold.

% PROVENANCE-BEGIN abstract-compact-inversion
% note: Compact operands: build/results/manuscript_results.csv; block=abstract-compact-inversion
% source: build/results/domains/raw/ideology_k_gap_accounting_both_universes.csv
% row:
%   key: election=2022; ideological_universe=seat_winning; k=0; coalition_id=PP|REPUBLICANOS|PSC|UNIÃO|PATRIOTA|NOVO|PL
%   row_at_generation: 9333
%   fields:
%     party_count: 7
%     vote_share: 0.4526056142903306
%     seats: 258
%   display:
%     party_count: seven
%     vote_share: 45.26
%     seats: 258
% source: build/results/cabinet_sets/summary.csv
% row:
%   key: election_year=2014
%   row_at_generation: 1
%   fields:
%     inverted_sets: 1
% source: build/results/cabinet_sets/summary.csv
% row:
%   key: election_year=2018
%   row_at_generation: 2
%   fields:
%     inverted_sets: 0
% source: build/results/cabinet_sets/summary.csv
% row:
%   key: election_year=2022
%   row_at_generation: 3
%   fields:
%     inverted_sets: 1
% PROVENANCE-END abstract-compact-inversion

\end{abstract}

\section{Introduction}

Under proportional representation, vote--seat translation is usually assessed
party by party, by comparing each party's vote and seat shares \parencite{kaminski26_propor_repres_primer}. In
fragmented legislatures, however, governing and legislative majorities are
assembled from several parties. Combining parties can make individually modest
deviations consequential. Positive and negative deviations may cancel or
reinforce one another, allowing a party set to cross the legislative majority
threshold before its combined vote share reaches one half. I refer to this
configuration as a coalition inversion: a party set holds a seat majority while
receiving less than half of all valid national legislative votes.
 


\textcite{kurrild-klitgaard13_elect_inver_coalit_propor_repres} documents empirical cases of coalition inversions\footnote{ I use
  \emph{coalition inversion} as shorthand for Miller's \emph{coalitionwise
    election inversion}.} in Danish government formation. \textcite{miller14_elect_inver_under_propor_repres} establishes the
more general result that inversions remain possible even when seats are
allocated in a single national constituency, without an explicit seat threshold,
and by a highly proportional formula. These works settle the existence question.
They do not show how coalition inversions appear when national party seat totals
are themselves assembled from many district elections.

Brazil provides such a setting. The Chamber of Deputies is not formed by
applying one proportional allocation to national party votes. Its seats are
allocated separately across the twenty-six states and the Federal District.
District magnitudes and electoral weights vary, electoral rules and the units
over which votes were pooled changed across the three elections, and party
support is distributed unevenly across the country. Several institutional and
geographical mediations therefore stand between a party's national vote share
and its final national seat share \parencite{bochsler2026winners}.  


Party-level deviations emerge from twenty-seven district-level translations
operating under heterogeneous conditions, rather than only from rounding a
single national allocation. The organizing question is therefore: when do these
deviations combine so that a politically meaningful set of parties holds a
legislative majority without a national vote majority? The contribution is to
show which governing combinations and ideologically structured majorities
benefit from these deviations, and how their advantages reflect coalition
composition, within-district allocation, and the distribution of
electoral support.

Coalition inversion becomes more politically informative once the combinations
of parties under consideration are given substantive structure. I examine the
phenomenon in two domains that capture distinct ways parties combine in
parliamentary politics. The first is realized cabinet composition. Changes in
ministerial participation generate successive cabinet party sets during a
legislative term; for each, I compare the Chamber seats held by its members with
their combined vote share in the election that produced the legislature.

The second domain examines inversions in the ideological structure of the party
system. Following \textcite{axelrod1970conflict}, I identify minimal connected
winning coalitions: contiguous blocs in the election-year left--right order that
hold a seat majority and lose it if narrowed. A one-gap sensitivity permits one
represented party inside the span to be omitted. Taken together, the two domains
locate coalition inversions in realized governing combinations and in
ideologically structured potential majorities. For both, I examine how
coalition-level departures from proportionality reflect vote-to-seat translation
within districts and differences in electoral weight across districts.

Inversions appear in both coalition domains, but their distribution differs
across elections. Two of the 35 distinct cabinet party sets invert: one observed for five days
under Dilma Rousseff following the 2014 election and one observed during the first 255
days of Lula III following the 2022 election. No cabinet party set linked to
the 2018 election inverts. Minimal connected
winning inversions, by contrast, appear in all three elections. The 2022
cases include a compact right-side bloc with 45.26 percent of the vote and
258 seats. Within-district and between-district components reinforce one
another in some inversions and offset one another in others.

% PROVENANCE-BEGIN introduction-cabinet-set-scope
% note: Compact operands: build/results/manuscript_results.csv; block=introduction-cabinet-set-scope
% source: build/results/cabinet_sets/summary.csv
% row:
%   key: election_year=all
%   row_at_generation: 4
%   fields:
%     distinct_sets: 35
%     inverted_sets: 2
% source: build/results/cabinet_sets/summary.csv
% row:
%   key: election_year=2018
%   row_at_generation: 2
%   fields:
%     inverted_sets: 0
% source: build/results/cabinet_sets/occurrences.csv
% row:
%   key: scenario_id=primary; occurrence_id=C2014-a264923b9881ebab3d9b@2016-04-14
%   row_at_generation: 5
%   fields:
%     days: 5
%     start_inclusive: 2016-04-14
%     end_exclusive: 2016-04-19
% source: build/results/cabinet_sets/occurrences.csv
% row:
%   key: scenario_id=primary; occurrence_id=C2022-259238bde0cf94feb03e@2023-01-01
%   row_at_generation: 41
%   fields:
%     days: 255
%     start_inclusive: 2023-01-01
%     end_exclusive: 2023-09-13
% PROVENANCE-END introduction-cabinet-set-scope


The remainder of the article is structured as follows. Section~\ref{sec:brazil-system} describes Brazil's electoral system;
Section~\ref{sec:mea} defines the measures and coalition domains;
Section~\ref{sec:data} presents the data; the remaining sections report and
discuss the results.


\section{Brazil's multidistrict open-list proportional system}
\label{sec:brazil-system}

Coalition majorities are formed from national party seat totals, but those
totals are not produced by a national allocation. In Brazil, elections to the
Chamber of Deputies take place separately in the twenty-six states and the
Federal District, each of which forms an electoral district. The 513 seats are
allocated within these twenty-seven districts and then summed nationally. The
Federal District is simply an eight-seat district, and no compensatory tier
subsequently adjusts the resulting party totals, unlike compensatory mixed
systems, where a second allocation stage can correct constituency-level
disproportionality \parencite{FlisGrofmanKaminski2025Introduction, salamon2026mixed}.

District magnitude varies considerably across states, from eight to seventy
seats. Eleven districts elect only eight deputies, the median magnitude is ten,
the mean nineteen, and S\~ao Paulo elects seventy. These differences imply
substantial variation in the effective conditions of representation across
districts: winning a seat is generally more demanding where fewer seats are
available. Because each district is allocated independently, the resulting
advantages and shortfalls carry into national party seat totals. 
% PROVENANCE-BEGIN district-magnitudes
% note: Compact operands: build/results/manuscript_results.csv; block=district-magnitudes
% source: build/results/accounting/tables/prose_analysis_summaries.csv
% row:
%   key: summary=eight-seat-districts; metric=count; aggregation=sum
%   row_at_generation: 24
%   fields:
%     value: 11
%   display:
%     value: Eleven
% source: build/results/accounting/tables/prose_analysis_summaries.csv
% row:
%   key: summary=district-magnitude; metric=S_d; aggregation=median
%   row_at_generation: 22
%   fields:
%     value: 10.0
%   display:
%     value: ten
% source: build/results/accounting/tables/prose_analysis_summaries.csv
% row:
%   key: summary=district-magnitude; metric=S_d; aggregation=mean
%   row_at_generation: 23
%   fields:
%     value: 19.0
%   display:
%     value: nineteen
% PROVENANCE-END district-magnitudes

Within each district, voters encounter an open-list election. They may enter the
identifying number of an individual candidate or cast a vote for a party label.
A candidate vote performs two functions at once. It contributes to the aggregate
vote of the party, electoral coalition, or federation through which the
candidate competes, and it raises that candidate's position within the relevant
open list. A party-label vote contributes to the first calculation but not to
the ranking among candidates. That is, votes first determine how many seats an
electoral list receives. Nominal votes then determine which candidates within
that list take the seats, subject to the individual eligibility requirements in
force for the election. A candidate can therefore benefit from votes cast for
copartisans or electoral partners, but cannot ordinarily occupy a seat unless
the list to which the candidate belongs wins one \parencite{nicolau2006openlist}.

Valid candidate and party-label votes are first aggregated for the electoral
unit in force in each district, whether a party, proportional coalition, or
federation. The electoral quotient is obtained from the total valid vote and the
number of seats available in the district, using the statutory rounding rule.
Each electoral unit then receives an initial allocation of seats according to
how many times its vote total reaches that quotient, subject to the
candidate-eligibility requirements in force for the election. Any seats left
unfilled are distributed one at a time by the Jefferson--D'Hondt
highest-averages procedure, with each unit's priority declining as it
accumulates seats. Eligibility to compete for these remaining seats varies
across election years, as described below \parencite{nicolauSchmitt1995}.

Eligibility rules changed across elections. In 2014, there was no fixed
candidate threshold, but only lists reaching the electoral quotient could
contest remainders. In 2018, candidates ordinarily needed 10 percent of the
quotient; all lists could contest remainders while an eligible candidate
remained, after which any still-unfilled seats were assigned by highest averages
without that nominal-vote minimum. In 2022, the 10 percent candidate threshold
remained for initial seats. For remainder seats, lists first needed 80 percent
of the quotient and candidates 20 percent; any seats still unfilled were then
assigned by highest averages among lists meeting the 80 percent requirement,
without the 20 percent candidate threshold.

Proportional electoral coalitions, or \emph{coligações proporcionais}, were
part of the 2014 and 2018 elections. They were formed separately within each
electoral district, so a party could enter different proportional arrangements
in different states. These were therefore not national alliances whose
composition automatically carried across all twenty-seven contests. Member
parties operated as a joint open list when votes were converted into seats:
candidate and party-label votes credited to any member entered the coalition's
district total, seats were assigned to the coalition as a whole, and the
highest nominal vote-getters across the combined list occupied them. A candidate
from one party could therefore be elected partly because votes cast for
candidates or labels of another coalition member raised the joint total
\parencite{nicolau2006openlist}.

This resembles list apparentement because parties pool votes before seats are
settled, but the Brazilian arrangement differs from the two-stage system studied
by \textcite{pukelsheimLeutgaeb2013apparentements}. There, alliance seats are
subsequently apportioned among component lists. In Brazil, a single joint open
list selected winners by candidate votes, so gains to particular members depended
on the distribution of nominal votes rather than on a fixed secondary
apportionment rule.


Constitutional reform in 2017 abolished proportional electoral coalitions
beginning with the 2020 elections. Party federations were introduced by
legislation in 2021 as a new form of multiparty organization for proportional
elections. Three federations competed in the 2022 general election:
PT--PCdoB--PV, PSDB--CIDADANIA, and PSOL--REDE. Like the former
\emph{coligações}, federations pooled their members' votes in a common open
list, so constituent-party seat totals resulted from competition within the
pooled list rather than from independent proportional allocations. Unlike the
former electoral coalitions, federations operate nationally and in principle
persist for at least four years. They have some common congressional procedures,
including leadership and committee-allocation arrangements, but do not merge
their constituent parties, which retain separate identities and organizational
autonomy \parencite{fischMesquita2022}.



Several mediations therefore stand between a party's national vote share and its national seat share. Within each district, finite magnitude, the quotient and highest-averages calculations, legal eligibility thresholds, and the number and relative strength of competing electoral units shape the conversion of votes into seats. Electoral pooling adds a further mediation by changing both the vote total on which seats are allocated and, through the common candidate list, the eventual distribution of those seats among member parties. These arrangements varied across districts in 2014 and 2018 and became nationally fixed through federations in 2022.

The national seat distribution then aggregates twenty-seven independently determined district outcomes. Chamber seats are apportioned by population, but the constitutional floor of eight and ceiling of seventy prevent exact population-to-seat equality across states \parencite{samuelsSnyder2001value}. Turnout and valid-vote rates also differ across elections and districts, so a seat can correspond to different numbers of valid votes in different states. Because parties distribute their support unevenly across the country, national vote shares can therefore be assembled from geographically different configurations with different representational consequences. Electoral geography interacts with district seat weights and local competitive conditions before the resulting party seats are summed nationally. The resulting disproportionality has multiple sources: allocation rules, district magnitude, malapportionment and electoral weight, electoral geography, and party competition constitute distinct channels through which party vote and seat shares can diverge \parencite{bochsler2026winners}.

This layered translation occurs in an unusually fragmented party environment. Measured by the Laakso--Taagepera index, the effective numbers of electoral parties were 14.11 in 2014, 18.02 in 2018, and 12.34 in 2022; the corresponding effective numbers of parliamentary parties were 13.33, 16.41, and 9.91 \parencite{laakso1979effective}. The exact number of allocation units was not identical to the number of party labels, especially when district coalitions or national federations pooled votes. Even so, each election produced a Chamber in which legislative majorities had to be assembled across numerous partisan units. 

% PROVENANCE-BEGIN party-fragmentation
% note: Compact operands: build/results/manuscript_results.csv; block=party-fragmentation
% source: build/results/accounting/tables/summaries/party_fragmentation_summary.csv
% row:
%   key: election_year=2014
%   row_at_generation: 1
%   fields:
%     effective_electoral: 14.108696059703476
%     effective_parliamentary: 13.332438320077005
%   display:
%     effective_electoral: 14.11
%     effective_parliamentary: 13.33
% source: build/results/accounting/tables/summaries/party_fragmentation_summary.csv
% row:
%   key: election_year=2018
%   row_at_generation: 2
%   fields:
%     effective_electoral: 18.015274717712103
%     effective_parliamentary: 16.408067834653032
%   display:
%     effective_electoral: 18.02
%     effective_parliamentary: 16.41
% source: build/results/accounting/tables/summaries/party_fragmentation_summary.csv
% row:
%   key: election_year=2022
%   row_at_generation: 3
%   fields:
%     effective_electoral: 12.33766514019305
%     effective_parliamentary: 9.908844459505252
%   display:
%     effective_electoral: 12.34
%     effective_parliamentary: 9.91
% PROVENANCE-END party-fragmentation

The national party seat totals produced by this process enter a presidential
system in which governing capacity ordinarily depends on multiparty
arrangements. Brazilian coalitional presidentialism links the presidency to
combinations of parties across cabinet and Congress: ministries are distributed
among several parties, while executive--legislative coordination is organized
through agreements among them
\parencite{abranches1988presidencialismo,figueiredo2000presidential}.
The Chamber assembled from twenty-seven district elections therefore supplies
the parliamentary resources from which governing coalitions are constructed.


The realized party seat totals incorporate the effects of rounding, highest
averages, district magnitude, eligibility requirements, electoral pooling,
apportionment, turnout, valid-vote variation, electoral geography, and
district-level competition. Comparing those totals with national proportional
quotas identifies the positive and negative representational deviations carried
by each party. That is, national proportionality provides a common baseline
against which to assess how Brazil’s layered vote-to-seat process generates
representational deviations and when those deviations become majority-changing
as parties are recombined into coalitions. If \(v_i\) is the national valid vote
total credited to party \(i\) and \(V\) is the corresponding total for all
parties, including those that win no seats, the party's reference quota is
\(513v_i/V\) seats. We use all valid federal-deputy votes, including votes cast
for parties that win no seats. Restricting the denominator to seat-winning
parties would instead measure proportionality conditional on electoral
inclusion. Here, exclusion remains part of the vote-to-seat transformation,
since votes excluded from representation can affect the allocation of seats
among represented parties
\parencite{kaminski18_spoil_effec_propor_repres_system}.

% ============================================================================
% NEW SECTION 3
% ============================================================================

\section{Coalition inversions, domains, and measurement}
\label{sec:mea}


Party seat shares need not match their national vote shares exactly. When
parties are combined, these departures from proportionality may offset or
reinforce one another. A coalition inversion occurs when their combined effect
is sufficient for a party set to reach the Chamber's 257-seat majority threshold
while still receiving less than half of the national vote.

To give this coalition-level phenomenon political content, I examine it in two
domains that capture distinct forms of party combination. Observed cabinets
identify realized governing combinations, showing when inversion appears in the
composition of government. Ideological connectedness identifies potential
parliamentary majorities structured by proximity in the party system. The two
domains therefore place the same representational phenomenon in realized
government composition and in the ideological structure of majority formation.

Cabinet coalitions are substantively meaningful because Brazilian presidents
ordinarily govern through multiparty arrangements. Cabinet participation records
a party combination that was not merely arithmetically possible but
institutionally realized in executive government. The relevant representational
question is, thus, whether the members of the party set holding ministries
collectively possessed a Chamber majority despite receiving less than half of
the valid federal-deputy vote in the election that produced that legislature.
Cabinet membership changes during a legislative term while the underlying
election result remains fixed. As parties enter or leave, the governing
combination can move into or out of inversion without any new electoral
allocation. The cabinet sequence therefore shows whether, and for how long,
realized governing sets hold a Chamber majority without a vote majority.
Observed cabinets cover only a small share of possible parliamentary majorities,
so I also examine potential coalitions structured by ideological proximity.
Following \textcite{axelrod1970conflict}, I use minimal connected winning coalitions: contiguous intervals
in the election-year ideological order that hold at least 257 seats and lose
that majority when narrowed.


\subsection{Measurement }

Let \(P\) denote the set of all parties receiving valid votes in a
federal-deputy election, including those that win no seats. For each party \(i\in P\),
let \(v_i\) be its national total of valid candidate and party-label
votes and \(s_i\) its Chamber seats. With \(V=\sum_{i\in P}v_i\) denoting
the total valid federal-deputy vote and \(S=513\), exact national
proportionality assigns party \(i\) the quota
\[
q_i=S\frac{v_i}{V}.
\]

The party seat differential is $d_i=s_i-q_i$; positive values denote
more seats than the national quota and negative values fewer. Because seats and
quotas each sum to $S$, party differentials sum to zero. For coalition
$C\subseteq P$,
\[
v_C=\sum_{i\in C}v_i,\quad s_C=\sum_{i\in C}s_i,\quad
q_C=S\frac{v_C}{V},\quad d_C=s_C-q_C=\sum_{i\in C}d_i.
\]

The coalition's departure from exact national proportionality is produced by
adding the positive and negative differentials of its member parties. A party
with a negative contribution can belong to an inverted coalition, provided that
the positive contributions elsewhere in the set more than offset it.

A coalition inversion occurs when the party set receives less than half of the
national valid federal-deputy vote but controls at least 257 Chamber seats:
\(\frac{v_C}{V}<\frac{1}{2}\) and \(s_C\geq257.\)
For a vote-minority coalition, the seat-majority condition is
$d_C\geq257-q_C$: its accumulated differential must bridge the distance from
its national proportional quota to 257 seats. Because $C$ and its complement
partition both votes and seats, an inversion simultaneously gives the complement
a vote majority and at most 256 seats. Throughout, $V$ includes all valid
federal-deputy votes, including votes for parties winning no seats.


Representation ratios express the same departures in relative rather than absolute terms \parencite{taagepera1980proportionality}. At the party and coalition levels,
% \[
% R_i=\frac{s_i/S}{v_i/V},
% \qquad
% R_C=\frac{s_C/S}{v_C/V}
% =\frac{s_C}{q_C}
% =1+\frac{d_C}{q_C}
% =\sum_{i\in C}\frac{q_i}{q_C}R_i.
% \]

\[
R_i=\frac{s_i/S}{v_i/V}=\frac{s_i}{q_i},
\qquad
R_C=\frac{s_C/S}{v_C/V}
=\frac{s_C}{q_C}
=1+\frac{d_C}{q_C}
=\sum_{i\in C}\frac{q_i}{q_C}R_i.
\]
A value of one denotes exact proportionality, a value above one denotes
overrepresentation, and a value below one denotes underrepresentation. The ratios
permit comparisons across parties and coalitions of different sizes. The
differentials remain the more direct quantities for the majority-threshold
question because they measure the relevant departure in seats.


The national coalition differential can be decomposed by adding and subtracting
the seats that exact proportionality would assign to the coalition separately
within each district:
\[
\begin{aligned}
d_C
&=
\left[
\sum_d
\left(
s_{Cd}-S_d\frac{v_{Cd}}{V_d}
\right)
\right]
+
\left[
\left(
\sum_d S_d\frac{v_{Cd}}{V_d}
\right)
-
\left(
S\frac{v_C}{V}
\right)
\right].
\end{aligned}
\]
The first term compares the coalition's actual seats in each district with the
number implied by its vote share in that district. I denote it the
within-district component, \(A_C\). The other component is the between-district
component, \(B_C\). Its first term sums the coalition's proportional quota in
each district. The second is the quota implied by the coalition's share of the
national valid vote. These quantities need not coincide because districts differ
in the number of seats they contribute relative to the number of valid votes
cast in them. Thus, \(B_C\) depends on the geographic distribution of coalition
support across districts with different electoral weights. Support concentrated
in districts whose share of Chamber seats exceeds their share of the national
valid vote raises the sum of district quotas relative to the national quota,
while support concentrated in districts with the opposite weighting lowers it.

Thus, \(d_C=A_C+B_C\). For comparisons across coalitions of different electoral
sizes, the two components can also be expressed relative to the coalition's
proportional quota: \(R_C-1=\frac{A_C}{q_C}+\frac{B_C}{q_C}.\) This expresses
the within- and between-district contributions to the coalition's relative
departure from proportionality.

In Brazil, \(A_C\) can reflect the mechanisms that operate within each district
when votes are converted into seats: finite district magnitude, the
electoral-quotient and highest-averages calculations, the year-specific
eligibility rules for parties and candidates, and the pooling of votes through
proportional electoral coalitions in 2014 and 2018 and party federations in
2022. These mechanisms determine how the coalition's actual district seats
depart from its district-level proportional quotas. By contrast, \(B_C\) arises
from differences in electoral weight across districts and the geographic
distribution of coalition support: district apportionment, turnout, and
valid-vote variation determine how many seats correspond to a given number of
votes in each district, while coalition geography determines whether its support
is concentrated in relatively more or less heavily weighted districts. Neither
component isolates the independent effect of any one mechanism; the
decomposition just locates the coalition differential as a within or between
districts phenomenon \parencite{bochsler2026winners}.


The primary ideological universe is the set of parliamentary parties,
\(P^+=\{i\in P:s_i>0\}.\) National proportionality is evaluated over \(P\),
while potential parliamentary coalitions are constructed over \(P^+\). Thus,
votes for parties without seats remain part of \(V\) and of party-level
proportionality accounting, but those parties cannot constitute or interrupt a
primary ideological coalition. 

Axelrod's connected-coalition criterion provides the baseline restriction
\parencite{axelrod1970conflict}. Let \(p_1,\ldots,p_n\) denote the existing election-year ideological order
restricted to \(P^+\), preserving the relative ordering of represented parties
without re-estimating their positions. A coalition \(C\subseteq P^+\) is connected when
its members form a contiguous interval in this filtered order. A connected
coalition is minimal winning when \(s_C\geq257\) and no proper connected subset
retains a Chamber majority. For an interval, removing either endpoint must
therefore leave fewer than 257 seats. Exact connectedness (\(k=0\)) provides the
baseline ideological domain. As a sensitivity analysis, I use a one-gap domain (\(k=1\)), which
allows at most one represented party inside a coalition's ideological
span to be omitted. A winning coalition is minimal within this domain
if no proper subset satisfying the same restriction retains a majority.\footnote{
Appendix~\ref{app:one-gap} gives the formal definition and summary results.
I also repeat both searches using the full election-year party order,
including parties that received votes but no seats, so that
non-seat-winning parties can interrupt connectedness. Connectedness and
minimality are recomputed in this universe, while the underlying vote-seat
data and national vote denominator remain unchanged.
Appendix~\ref{app:ideological-universes} compares the results.
Full coalition-level results are available in the replication outputs.}


\subsection{Data and operationalization}
\label{sec:data}

Vote and seat data for the 2014, 2018, and 2022 federal-deputy elections
come from the Tribunal Superior Eleitoral, accessed through the
\texttt{electionsBR} package \parencite{electionsbr}.\footnote{All 2022
results refer to the official seat vector in the frozen dataset, prior
to subsequent judicial retotalizations.} For each party, the national
vote total combines valid candidate votes and party-label votes across
the twenty-seven districts, while the seat total counts deputies elected
under its election-year label.\footnote{For 2014, the vote measure is
\texttt{QT\_VOTOS\_NOMINAIS + QT\_VOTOS\_LEGENDA}. For 2018 and 2022,
it is \texttt{QT\_VOTOS\_NOMINAIS\_VALIDOS +
QT\_TOTAL\_VOTOS\_LEG\_VALIDOS}.} This attribution also applies to
parties that competed through proportional electoral coalitions in
2014 and 2018 or federations in 2022. The resulting party totals connect
the district electoral outcomes to the party identities observed in
cabinet participation and ideological classifications.

Cabinet composition is reconstructed from publicly available online records for
the Dilma Rousseff, Michel Temer, Jair Bolsonaro, and Lula III administrations,
recording each officeholder's service dates and party affiliation at the time. I
map these dated affiliations to the party labels of the election that produced
the corresponding legislative term and treat each distinct party set within that
term as one cabinet configuration\footnote{For compact reference, cabinet
configurations are labeled \textit{YY-nn}, where \textit{YY} identifies the
election whose vote--seat distribution is used and \textit{nn} gives the
configuration's order of first appearance in the corresponding legislative-term
cabinet chronology. Thus, 14-05 denotes the fifth distinct configuration
associated with the 2014 election, 18-11 the eleventh associated with 2018, and
22-01 the first associated with 2022. These labels are identifiers, not dates;
a recurring configuration retains the same label.}.

Each cabinet configuration is evaluated against the vote and seat distribution
of the election that produced the corresponding Chamber: 2014 for the
2015--2018 term, 2018 for 2019--2022, and 2022 for Lula III through March
19, 2026. Parties that renamed or merged during a term are matched to the
election-year parties from which their vote and seat totals derive\footnote{
Brazilian parties are generally identified by their election-year party
abbreviations rather than by their full names.}. This yields
12 distinct cabinet configurations for the 2014 term, 18 for 2018, and 5 for
2022\footnote{Appendix~\ref{app:composition} reports their membership.}.

% PROVENANCE-BEGIN cabinet-set-construction
% note: Compact operands: build/results/manuscript_results.csv; block=cabinet-set-construction
% source: build/results/cabinet_sets/summary.csv
% row:
%   key: election_year=2014
%   row_at_generation: 1
%   fields:
%     distinct_sets: 12
%     analytical_periods: 20
%     observed_days: 1461
% source: build/results/cabinet_sets/summary.csv
% row:
%   key: election_year=2018
%   row_at_generation: 2
%   fields:
%     distinct_sets: 18
%     analytical_periods: 20
%     observed_days: 1461
% source: build/results/cabinet_sets/summary.csv
% row:
%   key: election_year=2022
%   row_at_generation: 3
%   fields:
%     distinct_sets: 5
%     analytical_periods: 13
%     observed_days: 1174
% PROVENANCE-END cabinet-set-construction

Cabinet affiliations are identified for 3,996 of the 4,096 days covered by the
chronology. The remaining 100 days involve unresolved affiliations for a small
number of officeholders. Alternative codings for the recorded cases do not
change which cabinet configurations invert or the 260 days for which an
inversion is observed.
% PROVENANCE-BEGIN cabinet-evidence-and-duration
% note: Compact operands: build/results/manuscript_results.csv; block=cabinet-evidence-and-duration
% source: build/results/cabinet_sets/summary.csv
% row:
%   key: election_year=all
%   row_at_generation: 4
%   fields:
%     observed_days: 4096
%     established_days: 3996
%     provisional_days: 100
%     inversion_days: 260
%     provisional_only_sets: 1
% note: V5 primary_assumptions.csv retains nine distinct unresolved people; sensitivity_summary.csv has 52 separate scenario calendars, each retaining 260 inversion days.
% PROVENANCE-END cabinet-evidence-and-duration

The ideological orders for 2014 and 2018 use the party classifications
of \textcite{bolognesi2023nova}. For 2022, I use the later 2018--2022
classification, drawing on the 2025 data release and associated article
\parencite{bolognesi2026desaparecimento}.
Each classification supplies the election-year party order used to construct
the coalition domains defined above.


\section{Results}
\label{sec:party}

Before turning to coalitions, I first examine the party-level seat differentials
\(d_i\) and representation ratios \(R_i\) produced in each election, which
describe the absolute and relative departures from national proportionality
that are subsequently combined at the coalition level.


In the 2014 election PMDB received 8.13 seats above its national proportional
quota, the largest positive differential that year. PSD and PTB received another
4.55 and 4.37 seats above their quotas, respectively. These gains coexisted with
deficits of 4.43 seats for PSDB and 4.20 for PSOL. The connected inversions
identified below all combine PMDB and PSD with PSDB: their majority advantage
therefore does not require every member to be overrepresented.

% PROVENANCE-BEGIN party-differentials-2014
% note: Compact operands: build/results/manuscript_results.csv; block=party-differentials-2014
% source: build/results/accounting/raw/party_accounting_all_years.csv
% row:
%   key: election_year=2014; party=PMDB
%   row_at_generation: 8
%   fields:
%     d_i: 8.133380861621642
%   display:
%     d_i: 8.13
% source: build/results/accounting/raw/party_accounting_all_years.csv
% row:
%   key: election_year=2014; party=PSD
%   row_at_generation: 20
%   fields:
%     d_i: 4.552732200018501
%   display:
%     d_i: 4.55
% source: build/results/accounting/raw/party_accounting_all_years.csv
% row:
%   key: election_year=2014; party=PTB
%   row_at_generation: 28
%   fields:
%     d_i: 4.374724383417064
%   display:
%     d_i: 4.37
% source: build/results/accounting/raw/party_accounting_all_years.csv
% row:
%   key: election_year=2014; party=PSDB
%   row_at_generation: 21
%   fields:
%     d_i: -4.4303847839739765
%   display:
%     d_i: 4.43
% source: build/results/accounting/raw/party_accounting_all_years.csv
% row:
%   key: election_year=2014; party=PSOL
%   row_at_generation: 24
%   fields:
%     d_i: -4.1975024814762625
%   display:
%     d_i: 4.20
% PROVENANCE-END party-differentials-2014

In 2018, PP, PR, and MDB had positive differentials of 8.76, 5.72, and
5.60 seats, while PSL and NOVO fell 7.70 and 6.35 seats below quota. PSL had
the largest national vote total but the largest negative differential, showing
that party size affects the magnitude of representation without fixing its
sign. Whether a party is over- or underrepresented nationally depends on both
its vote-to-seat conversion within districts and the distribution of its votes
across districts with different electoral weights.

% PROVENANCE-BEGIN party-differentials-2018
% note: Compact operands: build/results/manuscript_results.csv; block=party-differentials-2018
% source: build/results/accounting/raw/party_accounting_all_years.csv
% row:
%   key: election_year=2018; party=PP
%   row_at_generation: 47
%   fields:
%     d_i: 8.763869279337671
%   display:
%     d_i: 8.76
% source: build/results/accounting/raw/party_accounting_all_years.csv
% row:
%   key: election_year=2018; party=PR
%   row_at_generation: 50
%   fields:
%     d_i: 5.724395499520222
%   display:
%     d_i: 5.72
% source: build/results/accounting/raw/party_accounting_all_years.csv
% row:
%   key: election_year=2018; party=MDB
%   row_at_generation: 36
%   fields:
%     d_i: 5.600328135814278
%   display:
%     d_i: 5.60
% source: build/results/accounting/raw/party_accounting_all_years.csv
% row:
%   key: election_year=2018; party=PSL
%   row_at_generation: 59
%   fields:
%     d_i: -7.697795432903889
%   display:
%     d_i: 7.70
% source: build/results/accounting/raw/party_accounting_all_years.csv
% row:
%   key: election_year=2018; party=NOVO
%   row_at_generation: 37
%   fields:
%     d_i: -6.34667631209294
%   display:
%     d_i: 6.35
% PROVENANCE-END party-differentials-2018

In 2022, positive differentials were concentrated among large parties: PL,
UNI\~AO, PT, and PP gained 13.78, 11.10, 6.94, and 6.24 seats above quota.
Because $d_i=q_i(R_i-1)$, electoral size scales the seat consequence of a given
relative departure from proportionality. Large parties lying on the positive
side therefore placed substantial surpluses inside many available coalition
combinations.

% PROVENANCE-BEGIN party-differentials-2022
% note: Compact operands: build/results/manuscript_results.csv; block=party-differentials-2022
% source: build/results/accounting/raw/party_accounting_all_years.csv
% row:
%   key: election_year=2022; party=PL
%   row_at_generation: 79
%   fields:
%     d_i: 13.779741154081268
%   display:
%     d_i: 13.78
% source: build/results/accounting/raw/party_accounting_all_years.csv
% row:
%   key: election_year=2022; party=UNIÃO
%   row_at_generation: 98
%   fields:
%     d_i: 11.103563583757866
%   display:
%     d_i: 11.10
% source: build/results/accounting/raw/party_accounting_all_years.csv
% row:
%   key: election_year=2022; party=PT
%   row_at_generation: 92
%   fields:
%     d_i: 6.937954845575375
%   display:
%     d_i: 6.94
% source: build/results/accounting/raw/party_accounting_all_years.csv
% row:
%   key: election_year=2022; party=PP
%   row_at_generation: 83
%   fields:
%     d_i: 6.242080657901948
%   display:
%     d_i: 6.24
% PROVENANCE-END party-differentials-2022

Figure~\ref{fig:vote-seat} shows greater dispersion in representation ratios
among small parties, while larger parties generally lie closer to \(R_i=1\).
The direction of the deviation does not follow a clear size pattern in 2014 or
2018. PT and PSDB were underrepresented in 2014 while PMDB, of comparable
size, was overrepresented; in 2018, the largest party, PSL, was
underrepresented while several smaller parties were overrepresented. The 2022
profile is different: every party receiving more than 5 percent of the national
vote was overrepresented, while every party between 3 and 5 percent was
underrepresented. A clear association between larger size and
overrepresentation therefore appears only in 2022.

\begin{figure}[htbp]
\centering
\includegraphics[width=\textwidth]{../../../build/assets/party_representation_profile.pdf}
\caption{Party-level representation profiles, 2014, 2018, and 2022}
\label{fig:vote-seat}
\begin{flushleft}
\footnotesize Notes: Each point is a party. \(R_i=1\) denotes exact
proportionality. 
\end{flushleft}
\end{figure}

Representation ratios alone also do not identify the largest contributors to a
coalition surplus because $R_C$ is quota-weighted and $d_C=\sum_{i\in C}d_i$.
In 2022, PV had the higher ratio ($R_i=1.341$) but only a 1.52-seat surplus,
whereas PL's lower ratio of 1.162 produced 13.78 surplus seats. Majority status
therefore depends on the size-weighted combination of member gains and losses
and on whether their total bridges the distance to 257 seats.

% PROVENANCE-BEGIN party-ratios-2022
% note: Compact operands: build/results/manuscript_results.csv; block=party-ratios-2022
% source: build/results/accounting/raw/party_accounting_all_years.csv
% row:
%   key: election_year=2022; party=PV
%   row_at_generation: 94
%   fields:
%     R_i: 1.3405820708698333
%     d_i: 1.5243322058552404
%   display:
%     R_i: 1.341
%     d_i: 1.52
% source: build/results/accounting/raw/party_accounting_all_years.csv
% row:
%   key: election_year=2022; party=PL
%   row_at_generation: 79
%   fields:
%     R_i: 1.1616956031428576
%     d_i: 13.779741154081268
%   display:
%     R_i: 1.162
%     d_i: 13.78
% PROVENANCE-END party-ratios-2022

\subsection{Cabinet coalition inversions}
\label{sec:cabinets}

Two of the 35 distinct cabinet party sets combine a Chamber majority with less
than half of the national federal-deputy vote. They occur under Dilma Rousseff
and Lula III, following the 2014 and 2022 elections, respectively. No set
observed under Michel Temer or Jair Bolsonaro inverts. The two inversions differ
in their distance from a vote majority, the size of their seat advantage, and
their observed duration. Table~\ref{tab:cabinet-inversions} reports these sets and the within-district
and between-district components of their differentials.

% PROVENANCE-BEGIN cabinet-inversion-overview
% note: Compact operands: build/results/manuscript_results.csv; block=cabinet-inversion-overview
% source: build/results/cabinet_sets/summary.csv
% row:
%   key: election_year=all
%   row_at_generation: 4
%   fields:
%     distinct_sets: 35
%     inverted_sets: 2
%   display:
%     distinct_sets: 35
%     inverted_sets: Two
% source: build/results/cabinet_sets/summary.csv
% row:
%   key: election_year=2018
%   row_at_generation: 2
%   fields:
%     inverted_sets: 0
% source: build/results/accounting/raw/cabinet_party_set_accounting.csv
% row:
%   key: election_year=2014; cabinet_party_set_id=C2014-a264923b9881ebab3d9b
%   row_at_generation: 5
%   fields:
%     canonical_membership: PCdoB;PDT;PMDB;PR;PSD;PT;PTB
%     display_label: 14-05
%     administrations: dilma_2
% source: build/results/accounting/raw/cabinet_party_set_accounting.csv
% row:
%   key: election_year=2022; cabinet_party_set_id=C2022-259238bde0cf94feb03e
%   row_at_generation: 31
%   fields:
%     canonical_membership: MDB;PCdoB;PDT;PSB;PSD;PSOL;PT;REDE;UNIÃO
%     display_label: 22-01
%     administrations: lula_3
% PROVENANCE-END cabinet-inversion-overview

% Administration check: filter cabinet_analysis_periods.csv on administration=temer
% or administration=bolsonaro; neither subset has inversion_status=True.

\begin{table}[htbp]
\centering
\caption{Inverted cabinet party sets and their district components}
\label{tab:cabinet-inversions}
\footnotesize
\setlength{\tabcolsep}{2pt}
\renewcommand{\arraystretch}{1.06}
\input{../../../build/assets/table_observed_inversion_decomposition.tex}
\end{table}


The first inversion emerges during the contraction of Dilma Rousseff's
cabinet in April 2016. PRB ceases to be represented on March 18, and PP's
departure on April 14 leaves a party set with 259 seats and 46.54 percent
of the 2014 vote. Exact national proportionality would assign those
parties 238.73 seats, leaving them 18.27 seats short of a Chamber majority.
Their 20.27-seat differential, corresponding to \(R_C=1.085\), is
therefore sufficient to preserve the majority. Set 14-05 is observed for five days, from April 14 through April 18. PSD's departure
on April 19 reduces the party set to 223 seats and ends the inversion.

% PROVENANCE-BEGIN cabinet-dilma-2016
% note: Compact operands: build/results/manuscript_results.csv; block=cabinet-dilma-2016
% source: build/results/accounting/raw/cabinet_party_set_accounting.csv
% row:
%   key: election_year=2014; cabinet_party_set_id=C2014-a264923b9881ebab3d9b
%   row_at_generation: 5
%   fields:
%     canonical_membership: PCdoB;PDT;PMDB;PR;PSD;PT;PTB
%     display_label: 14-05
%     votes: 45305141
%     seats: 259
%     vote_share: 0.4653584948188879
%     q_C: 238.7289078420895
%     q_C_exact: 1223238807//5123966
%     d_C: 20.271092157910495
%     R_C: 1.0849125995722273
%   display:
%     display_label: 14-05
%     seats: 259
%     vote_share: 46.54
%     q_C: 238.73
%     d_C: 20.27
%     R_C: 1.085
% source: build/results/cabinet_sets/occurrences.csv
% row:
%   key: scenario_id=primary; occurrence_id=C2014-a264923b9881ebab3d9b@2016-04-14
%   row_at_generation: 5
%   fields:
%     start_inclusive: 2016-04-14
%     end_exclusive: 2016-04-19
%     days: 5
%   display:
%     days: five
% source: build/results/accounting/raw/cabinet_party_set_accounting.csv
% row:
%   key: election_year=2014; cabinet_party_set_id=C2014-8d3ff5ab64ec71cc91c5
%   row_at_generation: 6
%   fields:
%     canonical_membership: PCdoB;PDT;PMDB;PR;PT;PTB
%     seats: 223
%   display:
%     seats: 223
% note: Proportional majority shortfall is 257-q_C=18.271092157910495. Occurrence ends are exclusive. V5 transitions preserve the historical entry/exit narrative.
% PROVENANCE-END cabinet-dilma-2016

% Threshold calculation from CV5-005: 257 - q_C = 18.271092157910495;
% display as 18.27 seats, without introducing a separate threshold symbol.
% Membership-date source: data/cabinet/cabinet_periods.csv (CV5-002 through CV5-006).
% Stable REPUBLICANOS ceases representation on 2016-03-18;
% PP leaves on 2016-04-14; PSD leaves on 2016-04-19.
% Use PRB in the manuscript for the corresponding 2014 party label.

The Temer sequence does not produce an inversion. The cabinet party set
observed from April 2 through April 5, 2018 comes close, with 268 seats
and 50.03 percent of the 2014 vote, but it still has a vote majority.
From April 6, the remaining parties have 247 seats and 45.83 percent of
the vote. The contraction thus moves from a majority in both measures
to a minority in both, rather than leaving a vote-minority seat majority.
A positive seat differential is not, by itself, evidence of an inversion.
% Retire cabinet-temer-2017: its old inversion case_id no longer exists.
% The new non-inversion examples are selected by dates and membership,
% not by treating the reused 2017.1 label as the same observation.

% PROVENANCE-BEGIN cabinet-temer-no-inversion
% note: Compact operands: build/results/manuscript_results.csv; block=cabinet-temer-no-inversion
% source: build/results/accounting/raw/cabinet_party_set_accounting.csv
% row:
%   key: election_year=2014; cabinet_party_set_id=C2014-b99061a13a4cbf708682
%   row_at_generation: 10
%   fields:
%     canonical_membership: DEM;PMDB;PP;PRB;PSD;PSDB;PTB;PV
%     seats: 268
%     vote_share: 0.5003050576961592
%   display:
%     seats: 268
%     vote_share: 50.03
% source: build/results/cabinet_sets/occurrences.csv
% row:
%   key: scenario_id=primary; occurrence_id=C2014-b99061a13a4cbf708682@2018-04-02
%   row_at_generation: 18
%   fields:
%     start_inclusive: 2018-04-02
%     end_exclusive: 2018-04-06
%     days: 4
% source: build/results/accounting/raw/cabinet_party_set_accounting.csv
% row:
%   key: election_year=2014; cabinet_party_set_id=C2014-3a21551efde30c618990
%   row_at_generation: 11
%   fields:
%     canonical_membership: PMDB;PP;PRB;PSD;PSDB;PTB;PV
%     seats: 247
%     vote_share: 0.4583403702686963
%   display:
%     seats: 247
%     vote_share: 45.83
% source: build/results/cabinet_sets/occurrences.csv
% row:
%   key: scenario_id=primary; occurrence_id=C2014-3a21551efde30c618990@2018-04-06
%   row_at_generation: 19
%   fields:
%     start_inclusive: 2018-04-06
%     end_exclusive: 2018-07-05
% PROVENANCE-END cabinet-temer-no-inversion

The cabinet sequence linked to 2018 also contains no inversion. Set 18-11
holds 222 seats with 39.79 percent of the vote. Its 17.86-seat surplus over
the national proportional quota increases the parliamentary weight of its
members, but leaves their combination 35 seats short of a Chamber majority.
Overrepresentation therefore strengthens this cabinet party set without
supplying a majority.
% Threshold calculation from CV5-033: 257 - q_C = 52.86004690009656;
% display as 52.86 seats.
% Retire the old cabinet/2018/2021.3/2022.1 inversion key and its
% 250-seat successor. Neither identifies the new observations.

% PROVENANCE-BEGIN cabinet-bolsonaro-no-inversion
% note: Compact operands: build/results/manuscript_results.csv; block=cabinet-bolsonaro-no-inversion
% source: build/results/cabinet_sets/summary.csv
% row:
%   key: election_year=2018
%   row_at_generation: 2
%   fields:
%     inverted_sets: 0
% source: build/results/accounting/raw/cabinet_party_set_accounting.csv
% row:
%   key: election_year=2018; cabinet_party_set_id=C2018-a86e04ffb7898c05a4c4
%   row_at_generation: 23
%   fields:
%     canonical_membership: DEM;PP;PR;PRB;PSC;PSD;PSL
%     display_label: 18-11
%     seats: 222
%     vote_share: 0.3979336317736909
%     q_C: 204.13995309990344
%     d_C: 17.860046900096567
%     R_C: 1.0874892279972068
%   display:
%     display_label: 18-11
%     seats: 222
%     vote_share: 39.79
%     q_C: 204.14
%     d_C: 17.86
%     R_C: 1.087
% source: build/results/cabinet_sets/occurrences.csv
% row:
%   key: scenario_id=primary; occurrence_id=C2018-a86e04ffb7898c05a4c4@2021-07-28
%   row_at_generation: 33
%   fields:
%     start_inclusive: 2021-07-28
%     end_exclusive: 2022-03-24
%     days: 239
% note: 52.86 is the display of 257-q_C; no electoral quantity changes over this occurrence.
% PROVENANCE-END cabinet-bolsonaro-no-inversion


The second cabinet inversion appears at the beginning of Lula III. Its members
hold 263 seats with 48.89 percent of the 2022 vote. Their proportional quota is
250.80 seats, so the 12.20-seat differential (\(R_C=1.049\)) is sufficient to
bridge the 6.20-seat gap to a majority. The coalition therefore holds a larger
seat majority than the Dilma inversion despite a smaller differential, because
it begins much closer to the 257-seat threshold. This configuration persists
for 255 days, until PP and REPUBLICANOS enter the cabinet on September 13,
2023. Their entry raises the party set to 350 seats and 63.79 percent of the
vote, ending the inversion by adding a vote majority rather than by removing
the seat majority, as occurred in the Dilma case.

% PROVENANCE-BEGIN cabinet-lula-transition
% note: Compact operands: build/results/manuscript_results.csv; block=cabinet-lula-transition
% note: Previous display assertion display_label: 22-01 (no longer printed in this paragraph).
% source: build/results/accounting/raw/cabinet_party_set_accounting.csv
% row:
%   key: election_year=2022; cabinet_party_set_id=C2022-259238bde0cf94feb03e
%   row_at_generation: 31
%   fields:
%     canonical_membership: MDB;PCdoB;PDT;PSB;PSD;PSOL;PT;REDE;UNIÃO
%     display_label: 22-01
%     seats: 263
%     vote_share: 0.48888043147286714
%     q_C: 250.79566134558084
%     d_C: 12.20433865441916
%     R_C: 1.0486624792029489
%   display:
%     seats: 263
%     vote_share: 48.89
%     q_C: 250.80
%     d_C: 12.20
%     R_C: 1.049
% source: build/results/cabinet_sets/occurrences.csv
% row:
%   key: scenario_id=primary; occurrence_id=C2022-259238bde0cf94feb03e@2023-01-01
%   row_at_generation: 41
%   fields:
%     start_inclusive: 2023-01-01
%     end_exclusive: 2023-09-13
%     days: 255
%   display:
%     days: 255
% source: build/results/accounting/raw/cabinet_party_set_accounting.csv
% row:
%   key: election_year=2022; cabinet_party_set_id=C2022-d7aa8241899cd03a076c
%   row_at_generation: 32
%   fields:
%     canonical_membership: MDB;PCdoB;PDT;PP;PSB;PSD;PSOL;PT;REDE;REPUBLICANOS;UNIÃO
%     seats: 350
%     vote_share: 0.6378913927154223
%   display:
%     seats: 350
%     vote_share: 63.79
% note: The proportional majority shortfall is 257-q_C, displayed as 6.20. The entry on September 13 starts the first occurrence of 22-02; later recurrences do not extend that first interval.
% PROVENANCE-END cabinet-lula-transition

% Threshold calculation from CV5-041: 257 - q_C = 6.20433865441916;
% display as 6.20 seats. The successor begins on 2023-09-13.


The positive within-district component of the two inversions extends to all 35
cabinet configurations: \(A_C>0\) in every case, while \(B_C>0\) in only
thirteen. Cabinet composition helps explain this regularity. In every
configuration, the mean national vote share of cabinet-member parties exceeds
that of nonmember parties in the same election-year party system, with the
difference ranging from 1.46 to 4.95 percentage points. Larger parties are also
usually on the positive side of the within-district component: among parties
receiving at least 5 percent of the national vote, \(A_i>0\) for six of seven in
2014, eight of nine in 2018, and all seven in 2022\footnote{
Here \(A_i\) denotes the within-district component for party \(i\).
Appendix~\ref{app:party-contributions} gives the party-to-coalition
identities and selected member-party contributions.}. In 34 of the 35 cabinet
configurations, the positive \(A_i\) contributions of members above this
benchmark exceed the negative contributions of all members combined; only one
configuration also requires positive contributions from smaller parties to keep
\(A_C>0\). Cabinet composition therefore repeatedly brings together relatively
large parties whose within-district contributions are predominantly positive.
No comparable sign pattern appears for \(B_C\).

The two inverted cabinet party sets differ in how the positive \(A_C\)
interacts with between-district weighting. For 14-05, \(A_C=15.74\)
raises the coalition's quota of 238.73 to \(q_C+A_C=254.47\) seats,
still short of a Chamber majority. Adding \(B_C=4.53\) brings the
observed total to 259 seats. For 22-01, by contrast, \(A_C=13.20\)
already exceeds the 6.20 seats separating the quota from a majority:
\(q_C+A_C=263.99\). The negative between-district component,
\(B_C=-0.99\), reduces the total to 263 seats without removing the
majority. Thus, within-district gains characterize both inversions, while
between-district weighting reinforces the Dilma coalition and offsets the
Lula coalition. Only in the former is the positive between-district
component also needed to bridge the proportional distance to 257 seats.

% PROVENANCE-BEGIN cabinet-decomposition-thresholds
% note: Compact operands: build/results/manuscript_results.csv; block=cabinet-decomposition-thresholds
% source: build/results/accounting/raw/cabinet_party_set_accounting.csv
% row:
%   key: election_year=2014; cabinet_party_set_id=C2014-a264923b9881ebab3d9b
%   row_at_generation: 5
%   fields:
%     canonical_membership: PCdoB;PDT;PMDB;PR;PSD;PT;PTB
%     display_label: 14-05
%     q_C: 238.7289078420895
%     A_C: 15.742022264747815
%     B_C: 4.529069893162678
%     seats: 259
%   display:
%     A_C: 15.74
%     q_C: 238.73
%     B_C: 4.53
%     seats: 259
% source: build/results/accounting/raw/cabinet_party_set_accounting.csv
% row:
%   key: election_year=2022; cabinet_party_set_id=C2022-259238bde0cf94feb03e
%   row_at_generation: 31
%   fields:
%     canonical_membership: MDB;PCdoB;PDT;PSB;PSD;PSOL;PT;REDE;UNIÃO
%     display_label: 22-01
%     q_C: 250.79566134558084
%     A_C: 13.196247377886202
%     B_C: -0.9919087234670415
%     seats: 263
%   display:
%     A_C: 13.20
%     B_C: -0.99
%     seats: 263
% note: Within-adjusted quotas are q_C+A_C, computed without rounding and displayed as 254.47 and 263.99; proportional shortfalls are 257-q_C.
% PROVENANCE-END cabinet-decomposition-thresholds

% Derived from CV5-005: q_C + A_C = 254.47093010683733; display 254.47.
% Derived from CV5-041: q_C + A_C = 263.991908723467; display 263.99.
% Derived from CV5-041: 257 - q_C = 6.20433865441916; display 6.20.
% Retire the obsolete Dilma within-adjusted-quota summary and the old
% Temer/Bolsonaro inversion references; select the surviving rows by
% analytical_period_id and verify their dates and election-year members.


Overrepresentation is much more common than inversion: 33 of the 35 cabinet sets
have $R_C>1$, while only two invert. As Figure~\ref{fig:timeline} shows, most overrepresented
sets remain either below both majority thresholds or above both.

% PROVENANCE-BEGIN cabinet-overrepresentation
% note: Compact operands: build/results/manuscript_results.csv; block=cabinet-overrepresentation
% source: build/results/cabinet_sets/summary.csv
% row:
%   key: election_year=all
%   row_at_generation: 4
%   fields:
%     distinct_sets: 35
%     inverted_sets: 2
%     overrepresented_sets: 33
%     underrepresented_sets: 2
%   display:
%     distinct_sets: 35
%     overrepresented_sets: 33
%     inverted_sets: two
%     underrepresented_sets: two
% source: build/results/accounting/raw/cabinet_party_set_accounting.csv
% row:
%   key: election_year=2018; cabinet_party_set_id=C2018-72d588ebc21b8a1a1f9e
%   row_at_generation: 15
%   fields:
%     canonical_membership: DEM;NOVO;PP;PSDB;PSL
%     display_label: 18-03
%     analytical_period_labels: 2020.2;2020.4
%     occurrence_count: 2
% source: build/results/accounting/raw/cabinet_party_set_accounting.csv
% row:
%   key: election_year=2018; cabinet_party_set_id=C2018-3d6dc0b0e630f95b3cc9
%   row_at_generation: 30
%   fields:
%     canonical_membership: DEM;NOVO;PRB;PSD;PSL
%     display_label: 18-18
%     analytical_period_labels: 2022.7
%     evidence_status: provisional_only
%     total_observed_days: 2
%     provisional_days: 2
% PROVENANCE-END cabinet-overrepresentation


\begin{figure}[htbp]
\centering
\includegraphics[width=.98\textwidth]{../../../build/assets/observed_coalition_timeline.pdf}
\caption{Cabinet party-set vote shares, seat shares, and representation ratios}
\label{fig:timeline}

\end{figure}

Inversions occupy 5 of 1,461 covered days after 2014, none of 1,461 after 2018, and 255 of 1,174 after 2022, for 260 of 4,096
days in total.\footnote{
Appendix~\ref{app:composition} reports covered days by cabinet
configuration. Figure~\ref{fig:cabinet-membership-inversions} in
Appendix~\ref{app:cabinet-interval-bridge} gives the dated intervals
of the two cabinet inversions.}
The episodes end differently: the Dilma set loses its seat
majority, while the Lula set acquires a vote majority. A fixed electoral
distribution can therefore sustain different majority relations as cabinet
membership changes.

% PROVENANCE-BEGIN cabinet-inversion-duration
% note: Compact operands: build/results/manuscript_results.csv; block=cabinet-inversion-duration
% note: Previous display assertion inversion_episodes: two (no longer printed in this paragraph).
% source: build/results/cabinet_sets/summary.csv
% row:
%   key: election_year=2014
%   row_at_generation: 1
%   fields:
%     observed_days: 1461
%     inversion_days: 5
%   display:
%     observed_days: 1,461
%     inversion_days: 5
% source: build/results/cabinet_sets/summary.csv
% row:
%   key: election_year=2018
%   row_at_generation: 2
%   fields:
%     observed_days: 1461
%     inversion_days: 0
% source: build/results/cabinet_sets/summary.csv
% row:
%   key: election_year=2022
%   row_at_generation: 3
%   fields:
%     observed_days: 1174
%     inversion_days: 255
%   display:
%     observed_days: 1,174
%     inversion_days: 255
% source: build/results/cabinet_sets/summary.csv
% row:
%   key: election_year=all
%   row_at_generation: 4
%   fields:
%     observed_days: 4096
%     inversion_days: 260
%     inversion_episodes: 2
%   display:
%     observed_days: 4,096
%     inversion_days: 260
% PROVENANCE-END cabinet-inversion-duration


% Full-calendar totals: inversion_days = 5 + 0 + 255 = 260;
% covered_days = 1461 + 1461 + 1174 = 4096;
% established_days = 1383 + 1439 + 1174 = 3996;
% provisional_days = 78 + 22 + 0 = 100.
% Both inversion rows in cabinet_analysis_periods.csv have provisional_days=0.
% Do not replace covered_days with established_days when reporting these shares.

\subsection{Ideologically constrained potential coalitions}
\label{sec:ideology}

There are minimal connected winning inversions in all three elections:
four of eight minimal connected winning coalitions invert in 2014, one
of eight in 2018, and two of four in 2022\footnote{
Appendix~\ref{app:minimal-connected-winning} reports all twenty
baseline minimal connected winning intervals, including those
with vote majorities.}. These coalitions are
constructed over 28, 30, and 23 seat-winning parties, respectively.
The counts establish that inversions occur in each election; the figures
below show where those cases lie in the ideological ordering and how
they relate to the wider set of connected majorities. Figure~\ref{fig:minimal-connected-inversions}
identifies the minimal inversion intervals themselves, while
Figure~\ref{fig:heatmaps} places them within the full exact-connected
interval space.
% PROVENANCE-BEGIN ideological-k0-summary
% note: Compact operands: build/results/manuscript_results.csv; block=ideological-k0-summary
% source: build/results/domains/tables/ideological_universe_comparison.csv
% row:
%   key: election=2014; ideological_universe=seat_winning; k=0
%   row_at_generation: 1
%   fields:
%     minimal_inversions: 4
%     minimal_seat_majority_coalitions: 8
%     ideological_party_count: 28
%   display:
%     minimal_inversions: four
%     minimal_seat_majority_coalitions: eight
%     ideological_party_count: 28
% source: build/results/domains/tables/ideological_universe_comparison.csv
% row:
%   key: election=2018; ideological_universe=seat_winning; k=0
%   row_at_generation: 5
%   fields:
%     minimal_inversions: 1
%     minimal_seat_majority_coalitions: 8
%     ideological_party_count: 30
%   display:
%     minimal_inversions: one
%     minimal_seat_majority_coalitions: eight
%     ideological_party_count: 30
% source: build/results/domains/tables/ideological_universe_comparison.csv
% row:
%   key: election=2022; ideological_universe=seat_winning; k=0
%   row_at_generation: 9
%   fields:
%     minimal_inversions: 2
%     minimal_seat_majority_coalitions: 4
%     ideological_party_count: 23
%   display:
%     minimal_inversions: two
%     minimal_seat_majority_coalitions: four
%     ideological_party_count: 23
% PROVENANCE-END ideological-k0-summary

Figure~\ref{fig:minimal-connected-inversions} marks PT, the principal organizational and electoral pole of the
Brazilian left, as a common reference point. This distinguishes connected
majorities extending through PT from those assembled entirely to its right. The
four 2014 connected inversions contain twelve to fifteen parties. PSB--PTN
begins on the center-left; the other three begin at parties classified as
center-right, and all extend into the right. None includes PT, despite PT
holding the presidency under Dilma Rousseff and anchoring the left side of the
party system. The connected inversion region instead begins to its right: every
case combines PMDB and PSD with PSDB. The strongest case, the thirteen-party
PTB--PR interval, reaches exactly 257 seats with 47.45 percent of all valid
votes. The 2014 pattern is therefore not a single anomalous coalition but a set
of overlapping combinations, organized around a common group of parties, that
can secure a Chamber majority without PT and without a national vote majority.

% PROVENANCE-BEGIN ideological-k0-2014
% note: Compact operands: build/results/manuscript_results.csv; block=ideological-k0-2014
% source: build/results/domains/tables/ideological_universe_comparison.csv
% row:
%   key: election=2014; ideological_universe=seat_winning; k=0
%   row_at_generation: 1
%   fields:
%     minimal_inversions: 4
%     inverted_members_min: 12
%     inverted_members_max: 15
%   display:
%     minimal_inversions: four
%     inverted_members_min: twelve
%     inverted_members_max: fifteen
% source: build/results/accounting/tables/prose_analysis_summaries.csv
% row:
%   key: summary=seat-winning-2014-k0-other-inversions; metric=count; aggregation=sum_minus_one
%   row_at_generation: 25
%   fields:
%     value: 3
%   display:
%     value: three
% source: build/results/domains/raw/ideology_k_gap_accounting_both_universes.csv
% row:
%   key: election=2014; ideological_universe=seat_winning; k=0; coalition_id=PTB|PT DO B|SOLIDARIEDADE|PMN|PHS|PMDB|PSD|PSDB|PTN|PRTB|PROS|PRP|PR
%   row_at_generation: 188
%   fields:
%     party_count: 13
%     coalition_label: PTB--PR
%     vote_share: 0.4744616510767348
%     seats: 257
%     q_C: 243.39882700236495
%     d_C: 13.601172997635034
%     R_C: 1.0558801912282958
%   display:
%     party_count: thirteen
%     coalition_label: PTB--PR
%     vote_share: 47.45
%     seats: 257
%     q_C: 243.40
%     d_C: 13.60
% PROVENANCE-END ideological-k0-2014

The 2018 configuration differs. PT--PSDB is the only exact-connected
minimal inversion, and its fourteen parties span the ideological order from
the principal party of the left through PSDB. They receive 49.95 percent of
all valid votes and hold 267 seats. Unlike the overlapping 2014 cases, the
2018 inversion therefore requires one broad connected combination cutting
across much of the party order, and its electoral support lies only marginally
below a vote majority. The other seven minimal connected majorities have vote
majorities. This potential inversion also contrasts with the  cabinet
domain, in which no party set linked to the 2018 election inverts.

% PROVENANCE-BEGIN ideological-k0-2018
% note: Compact operands: build/results/manuscript_results.csv; block=ideological-k0-2018
% source: build/results/domains/raw/ideology_k_gap_accounting_both_universes.csv
% row:
%   key: election=2018; ideological_universe=seat_winning; k=0; coalition_id=PT|PDT|PSB|REDE|PPS|PV|PTB|AVANTE|SOLIDARIEDADE|PMN|PHS|MDB|PSD|PSDB
%   row_at_generation: 4161
%   fields:
%     party_count: 14
%     vote_share: 0.4995104320302238
%     seats: 267
%     q_C: 256.24885163150481
%     d_C: 10.751148368495176
%     R_C: 1.0419558889729414
%   display:
%     party_count: fourteen
%     vote_share: 49.95
%     seats: 267
%     q_C: 256.25
%     d_C: 10.75
% source: build/results/domains/tables/ideological_universe_comparison.csv
% row:
%   key: election=2018; ideological_universe=seat_winning; k=0
%   row_at_generation: 5
%   fields:
%     non_inverted_minimal_majorities: 7
%   display:
%     non_inverted_minimal_majorities: seven
% source: build/results/cabinet_sets/summary.csv
% row:
%   key: election_year=2018
%   row_at_generation: 2
%   fields:
%     inverted_sets: 0
% PROVENANCE-END ideological-k0-2018


In 2022, the structure of available connected majorities differs across the
party order. PT--PP and PSB--PSC begin on the left and center-left and require
fifteen and sixteen parties, respectively, to assemble a majority. Both have
vote majorities even though their combined representation is below exact
national proportionality. The inverted majorities are assembled farther to
the right from substantially fewer parties: ten in MDB--UNI\~AO and seven in
PP--PL.

% PROVENANCE-BEGIN ideological-k0-2022-membership
% note: Compact operands: build/results/manuscript_results.csv; block=ideological-k0-2022-membership
% source: build/results/domains/raw/ideology_k_gap_accounting_both_universes.csv
% row:
%   key: election=2022; ideological_universe=seat_winning; k=0; coalition_id=PT|PSB|REDE|PDT|PV|SOLIDARIEDADE|CIDADANIA|AVANTE|MDB|PSDB|PSD|PODE|PROS|PTB|PP
%   row_at_generation: 9138
%   fields:
%     party_count: 15
%     d_C: -9.3576267657920269
%   display:
%     party_count: fifteen
%     d_C: 9.36
% source: build/results/domains/raw/ideology_k_gap_accounting_both_universes.csv
% row:
%   key: election=2022; ideological_universe=seat_winning; k=0; coalition_id=PSB|REDE|PDT|PV|SOLIDARIEDADE|CIDADANIA|AVANTE|MDB|PSDB|PSD|PODE|PROS|PTB|PP|REPUBLICANOS|PSC
%   row_at_generation: 9160
%   fields:
%     party_count: 16
%     d_C: -15.098170666459209
%   display:
%     party_count: sixteen
%     d_C: 15.10
% source: build/results/domains/raw/ideology_k_gap_accounting_both_universes.csv
% row:
%   key: election=2022; ideological_universe=seat_winning; k=0; coalition_id=MDB|PSDB|PSD|PODE|PROS|PTB|PP|REPUBLICANOS|PSC|UNIÃO
%   row_at_generation: 9273
%   fields:
%     party_count: 10
%   display:
%     party_count: ten
% source: build/results/domains/raw/ideology_k_gap_accounting_both_universes.csv
% row:
%   key: election=2022; ideological_universe=seat_winning; k=0; coalition_id=PP|REPUBLICANOS|PSC|UNIÃO|PATRIOTA|NOVO|PL
%   row_at_generation: 9333
%   fields:
%     party_count: 7
%   display:
%     party_count: seven
% PROVENANCE-END ideological-k0-2022-membership

MDB--UNI\~AO receives 49.22 percent of the vote and holds 265 seats, whereas
PP--PL reaches 258 seats with only 45.26 percent. The latter result is driven
largely by the substantial positive differentials of PL and UNI\~AO, which
more than offset underrepresented parties inside the interval. The 2022
pattern therefore places the strongest vote--seat reversal in a relatively
compact bloc on the right: parties accounting for substantially less than half
of the national vote can be combined into a Chamber majority without extending
the coalition toward the left or center-left.

% PROVENANCE-BEGIN ideological-k0-2022-quantities
% note: Compact operands: build/results/manuscript_results.csv; block=ideological-k0-2022-quantities
% source: build/results/domains/raw/ideology_k_gap_accounting_both_universes.csv
% row:
%   key: election=2022; ideological_universe=seat_winning; k=0; coalition_id=MDB|PSDB|PSD|PODE|PROS|PTB|PP|REPUBLICANOS|PSC|UNIÃO
%   row_at_generation: 9273
%   fields:
%     vote_share: 0.492189812614362
%     q_C: 252.49337387116771
%     r_C: 4.506626128832295
%     d_C: 12.506626128832282
%     R_C: 1.0495324924257761
%     seats: 265
%   display:
%     vote_share: 49.22
%     q_C: 252.49
%     r_C: 4.51
%     d_C: 12.51
%     seats: 265
% source: build/results/domains/raw/ideology_k_gap_accounting_both_universes.csv
% row:
%   key: election=2022; ideological_universe=seat_winning; k=0; coalition_id=PP|REPUBLICANOS|PSC|UNIÃO|PATRIOTA|NOVO|PL
%   row_at_generation: 9333
%   fields:
%     vote_share: 0.4526056142903306
%     q_C: 232.18668013093958
%     r_C: 24.81331986906042
%     d_C: 25.813319869060408
%     R_C: 1.1111748522977425
%     seats: 258
%   display:
%     vote_share: 45.26
%     q_C: 232.19
%     r_C: 24.81
%     d_C: 25.81
%     seats: 258
% source: build/results/accounting/raw/party_accounting_all_years.csv
% row:
%   key: election_year=2022; party=PL
%   row_at_generation: 79
%   fields:
%     d_i: 13.779741154081268
%   display:
%     d_i: 13.78
% source: build/results/accounting/raw/party_accounting_all_years.csv
% row:
%   key: election_year=2022; party=UNIÃO
%   row_at_generation: 98
%   fields:
%     d_i: 11.103563583757866
%   display:
%     d_i: 11.10
% source: build/results/accounting/tables/prose_analysis_summaries.csv
% row:
%   key: summary=seat-winning-k0-totals; metric=minimal_inversions; aggregation=sum
%   row_at_generation: 2
%   fields:
%     value: 7
% PROVENANCE-END ideological-k0-2022-quantities

\begin{figure}[htbp]
\centering
\includegraphics[width=.88\textwidth]
  {../../../build/assets/minimal_connected_winning_inversions_3x1_diamond.png}
  \caption{Minimal connected winning inversions by election year. Brackets
    identify coalition endpoints and report vote shares and seats; the diamond
    marks PT as a common reference point.}
\label{fig:minimal-connected-inversions}
\end{figure}

Figure~\ref{fig:minimal-connected-inversions} reveals three different
coalitional structures behind the same formal outcome. In 2014, inversion is
available through several overlapping combinations to PT's right. In 2018, it
appears only in one broad interval extending from PT through PSDB. In 2022,
it again occurs to PT's right, but now includes the compact seven-party
PP--PL interval. Coalition inversion is therefore not attached to a fixed
ideological location or coalition size. What changes across elections is which
regions of the party order contain combinations whose accumulated
representational advantages are sufficient to turn a vote minority into a
seat majority.


Figure~\ref{fig:heatmaps} shows whether these minimal cases remain inverted as
their connected intervals expand. Blue cells have both vote and seat majorities,
red cells are minimal connected inversions, and orange cells are larger
connected coalitions that remain inverted. Whereas Figure~\ref{fig:minimal-connected-inversions} identifies the
membership and endpoints of the minimal cases, the heatmaps show how those cases
sit within the surrounding space of connected coalitions.

The structure differs substantially across elections. In 2014, several
minimal inversions lie in overlapping center-to-right spans, and some are
accompanied by orange cells. The inversion is therefore not confined to the
minimal winning coalitions themselves: some larger connected intervals
also retain a seat majority while remaining below 50 percent of the national
vote. In 2018, the blue region of vote-and-seat majorities contracts and its
boundary shifts toward intervals extending farther to the right of the party
ordering. At the same time, PT--PSDB is the only inversion and there are
no orange cells. This is a distinct feature of the 2018 configuration. The
PT--PSDB interval has 49.95 percent of the vote and 267 seats, but any
larger exact-connected interval containing it also has a vote majority. The
single inversion is therefore isolated within the surrounding connected
coalition space rather than forming part of a larger set of inverted intervals.

The 2022 pattern differs again. The blue region is smaller, reflecting a lower
proportion of connected coalitions with both vote and seat majorities, while the
inversion cells are concentrated on the right. Unlike 2018, orange cells
reappear: some connected expansions of the minimal right-side inversions
continue to hold a Chamber majority without reaching 50 percent of the national
vote. Connected majorities extending through PT, by contrast, require broader
intervals and have both vote and seat majorities. The 2022 result is therefore
not merely the existence of two minimal inversions, but an asymmetry in where
connected vote-and-seat majorities and connected vote-minority seat majorities
occur within the party ordering.

Figure~\ref{fig:minimal-connected-inversions} and
Figure~\ref{fig:heatmaps} thus reveal different aspects of the same result:
the first identifies the coalitions at the minimal winning boundary, while
the second shows whether inversion remains present as the connected
interval is expanded. The right-side concentration in 2022 is consistent
with the party-level results above, where large right and center-right parties,
particularly PL and UNIÃO, carry substantial positive seat differentials
that can reinforce one another when adjacent parties are combined.

\begin{figure}[htbp]
\centering
\includegraphics[width=.31\textwidth]{../../../build/assets/ideological_interval_heatmap_2014.pdf}\hfill
\includegraphics[width=.31\textwidth]{../../../build/assets/ideological_interval_heatmap_2018.pdf}\hfill
\includegraphics[width=.31\textwidth]{../../../build/assets/ideological_interval_heatmap_2022.pdf}
\par\vspace{2pt}
\includegraphics[width=.95\textwidth]{../../../build/assets/ideological_interval_heatmap_legend.pdf}
\caption{Exact-connected (\(k=0\)) ideological interval status by election}
\label{fig:heatmaps}
\begin{flushleft}
\footnotesize
Note:
Start and end indices refer to the election-year ideological order of
seat-winning parties.
\end{flushleft}
\end{figure}

\begin{samepage}
Allowing one represented party inside the ideological span to be omitted
substantially expands the set of admissible minimal winning coalitions. Under
\(k=1\), 46 of 87 such coalitions invert in 2014, 42 of 108 in 2018, and 12
of 39 in 2022. Inversions therefore remain present in all three elections and
are not confined to the few exact-connected cases identified above. The
relaxation matters most substantively in 2018: whereas the sole
exact-connected inversion lies only marginally below a vote majority, the
strongest one-gap inversion, PCdoB--PODE with PSDB omitted, reaches 257 seats
with 47.59 percent of the national vote\footnote{
Appendix~\ref{app:one-gap} reports the formal definition, counts,
and strongest one-gap inversions; coalition-level results remain
in the replication outputs. Repeating the analysis in the full
election-year party order, including parties that received votes
but no seats, changes some memberships and removes the 2018
exact-connected inversion, while one-gap inversions remain in all
three elections. Appendix~\ref{app:ideological-universes} reports
the comparison.}.

% PROVENANCE-BEGIN ideological-robustness-overview
% note: Compact operands: build/results/manuscript_results.csv; block=ideological-robustness-overview
% source: build/results/domains/raw/ideology_k_gap_accounting_both_universes.csv
% row:
%   key: election=2018; ideological_universe=seat_winning; k=1; coalition_id=PCdoB|PT|PDT|PSB|REDE|PPS|PV|PTB|AVANTE|SOLIDARIEDADE|PMN|PHS|MDB|PSD|PODE
%   row_at_generation: 5105
%   fields:
%     coalition_label: PCdoB--PODE, omitting PSDB
%     vote_share: 0.4759116927539931
%     seats: 257
%   display:
%     coalition_label: PCdoB--PODE omitting PSDB
%     vote_share: 47.59
%     seats: 257
% PROVENANCE-END ideological-robustness-overview

\end{samepage}


\subsection{Comparing decomposition patterns across coalition domains}
\label{sec:decomposition}


The inversions identified in the two coalition domains share a positive
within-district component, but differ in how between-district weighting
contributes to their seat majorities. Table~\ref{tab:minimal-intervals}
reports the decomposition for the minimal connected inversions,
complementing the cabinet results in
Table~\ref{tab:cabinet-inversions}. Reading the tables together separates
two aspects: which component accounts for most of a coalition's seat
differential, and whether that component is sufficient to cover the
distance from its proportional quota to 257 seats. A component can account
for most of the differential without being sufficient to produce the
seat majority.

\begin{table}[htbp]
\centering
\caption{Minimal connected winning inversions and
seat-differential decomposition}
\label{tab:minimal-intervals}
\footnotesize
\setlength{\tabcolsep}{3pt}
\renewcommand{\arraystretch}{1.06}
\input{../../../build/assets/table_accounting_minimal_ideological.tex}
\end{table}
\FloatBarrier


The member-party contributions help explain the differences among the
coalition totals reported in the table.\footnote{Table~\ref{tab:coalition-party-component-extremes}
in Appendix~\ref{app:party-contributions} reports the largest positive and
negative \(A_i\) and \(B_i\) contributions for selected coalitions.}

MDB--UNIÃO provides the exceptional case within the minimal connected winning
inversions: it is the only case in which between-district weighting accounts for
more of the coalition's seat surplus than within-district allocation. Its
\(12.51\)-seat differential consists of \(A_C=2.95\) and \(B_C=9.55\). The
relatively small \(A_C\) results from substantial cancellation among its
members: positive \(A_i\) contributions from UNIÃO and MDB are largely offset by
negative contributions from PTB and PODE. Its positive \(B_C\), by contrast,
indicates that the coalition's support is distributed relatively toward
districts with greater electoral weight. Since its proportional quota of
\(252.49\) seats lies only \(4.51\) seats below a majority, the between-district
component alone is larger than that gap. MDB--UNIÃO is therefore the only
inversion whose seat differential is primarily associated with favorable
between-district weighting rather than within-district allocation.

PP--PL shows why the component that accounts for most of a coalition's
seat differential need not, by itself, span the distance from its proportional
quota to 257 seats. Its \(A_C=23.32\) accounts for most of its
\(25.81\)-seat differential, while \(B_C=2.49\) is comparatively small.
Yet the coalition's proportional quota is only \(232.19\) seats, leaving
it \(24.81\) seats short of a majority. The within-district component
therefore falls short of that distance on its own; adding the positive
between-district component brings the observed total to 258 seats. The
Dilma cabinet inversion and the 2014 PTB--PR interval have the same
 structure: \(A_C\) is the larger component, but both positive
components enter the differential that carries the coalition across 257.

The sensitivity analysis for ideological coalitions separates two features that
coincide in most of the cabinet and minimal-connected winning cases. Positive
within-district contributions remain pervasive: \(A_C>0\) in ninety-seven of the
one hundred minimal inversions under the one-gap relaxation. But within-district
allocation is not always the larger component. \(A_C>B_C\) in seventy-seven of
those one hundred cases, with substantial variation across elections: the
relationship holds in forty of forty-six inversions in 2014 and thirty-two of
forty-two in 2018, but only five of twelve in 2022\footnote{
  Appendix~\ref{app:one-gap} reports the component frequencies for the one-gap
  inversions.}. Thus, the predominance of within-district gains remains strong
overall, while the 2022 inversions reveal a different coalitional pattern in
which favorable between-district weighting often provides the larger
contribution to the seat surplus.

Restoring parties without seats to the ideological order reinforces the
need to distinguish these two patterns. MDB--UNIÃO then combines
\(A_C=-1.29\) with \(B_C=9.88\) and remains inverted\footnote{
Appendix~\ref{app:ideological-universes} compares the
seat-winning and all-party specifications.}. The decomposition therefore
shows both a persistent tendency for inverted coalitions to accumulate
within-district gains and meaningful variation in whether those gains or
between-district weighting account for the larger share of the resulting
seat differential.

% PROVENANCE-BEGIN all-party-mdb-uniao-components
% note: Compact operands: build/results/manuscript_results.csv; block=all-party-mdb-uniao-components
% source: build/results/domains/raw/ideology_k_gap_accounting_both_universes.csv
% row:
%   key: election=2022; ideological_universe=all_parties; k=0; coalition_id=MDB|PMN|PSDB|PSD|PMB|PODE|PROS|PRTB|AGIR|PTB|PP|DC|REPUBLICANOS|PSC|UNIÃO
%   row_at_generation: 25594
%   fields:
%     A_C: -1.2923532538854881
%     B_C: 9.8835518927226591
%   display:
%     A_C: -1.29
%     B_C: 9.88
% PROVENANCE-END all-party-mdb-uniao-components

\begin{figure}[htbp]
\centering
\includegraphics[width=\textwidth]{../../../build/assets/cross_domain_components.pdf}
\caption{Within-district and between-district contributions across coalition domains}
\label{fig:cross-domain-components}
\end{figure}
\FloatBarrier

Figure~\ref{fig:cross-domain-components} places the inversions within a
broader comparison of coalition representation. Panel A includes all
35 distinct election-year cabinet party sets, including those observed
under the provisional affiliation assumptions described in
Section~\ref{sec:data}. Panel B includes all twenty minimal connected
winning coalitions. Both panels express the components relative to each coalition's proportional
quota: the horizontal coordinate is \(100A_C/q_C\), and the vertical coordinate
is \(100B_C/q_C\). Since \(d_C=A_C+B_C\) and
\(R_C-1=d_C/q_C\), their sum is \(100(R_C-1)\). Coalitions with the
same representation ratio therefore lie on the same line with slope \(-1\),
while different values of \(R_C\) define parallel lines. The dashed diagonal
marks \(R_C=1\), or exact proportionality. 

All cabinet sets in Panel A lie to the right of \(A_C=0\), so their members
collectively benefit from the within-district translation of votes into seats.
What distinguishes them is how that support is distributed across districts with
different electoral weights. The 2014 election provides the clearest contrast.
All six Dilma configurations have \(B_C>0\): their member parties, taken
together, draw relatively more support from districts contributing more Chamber
seats per valid vote. All six Temer configurations have \(B_C<0\), indicating the opposite pattern of between-district weighting.. Because both sequences are evaluated against
the same 2014 electoral distribution, the sign reversal comes from cabinet
recomposition rather than from a new election. The transition from Dilma to
Temer therefore changes not only which parties compose the government, but also
how the government's electoral support is positioned within the territorial
weighting of the Chamber.

The Bolsonaro cabinet sets extend farther into the negative between-district
region. This reflects a recurrent territorial disadvantage in the geographic
distribution of their electoral support. PSL appears in fourteen of the eighteen
configurations.\footnote{Appendix~\ref{app:composition} lists the election-year
  member parties of each cabinet configuration.} All fourteen have \(B_C<0\),
and PSL contributes \(-5.28\) seats to the between-district total. NOVO, PRB,
PSDB, and PR also contribute negatively, whereas PP, MDB, PSC, PSD, and DEM
contribute positively. Taken together, Bolsonaro's cabinet parties therefore
tended to draw relatively more support from districts contributing fewer Chamber
seats per valid vote and relatively less support from districts with greater
electoral weight. The persistence of \(B_C<0\) across changing cabinet
configurations suggests that this was not specific to a single party set, but a
recurring feature of how the cabinet parties' combined support was distributed
across districts.

Unlike cabinet sets, minimal connected winning coalitions appear on both sides
of \(A_C=0\), reflecting differences in whether their members collectively gain
or lose seats in the within-district translation of votes into seats. Their
between-district contributions also vary. The 2018 PT--PSDB inversion, for
example, combines \(A_C=6.57\) with \(B_C=4.18\). Its members therefore benefit
both from the within-district allocation of seats and from the geographic
distribution of their support: taken together, they perform relatively better
in districts contributing more Chamber seats per valid vote. Most Bolsonaro
cabinet sets exhibit a different territorial profile. Their positive
within-district gains are partly offset because their members' combined support
is relatively concentrated in districts with lower electoral weight. The same
2018 electoral distribution can therefore advantage different party
combinations through different channels, depending on which parties are brought
together.

A broader difference between the two coalition domains is visible in their
distance from the \(R_C=1\) diagonal. The minimal connected winning coalitions
are generally much closer to exact proportionality than the cabinet party sets:
seventeen of the twenty lie between \(R_C=0.95\) and \(R_C=1.05\), compared with
fourteen of the thirty-five cabinet sets \footnote{
For the ideological coalitions, Appendix~\ref{app:minimal-connected-winning}
reports the vote shares and seat totals from which these ratios
can be recovered.}. Their combined seat shares therefore
tend to track their combined vote shares more closely than those of cabinet
party sets. Yet seven of these coalitions still invert. Coalition inversion
consequently need not arise from a large aggregate departure from
proportionality. For coalitions already close to the majority threshold, a
comparatively small net representational advantage can be sufficient to reverse
vote and seat majority status.

The 2022 cases also show the limit of this pattern. MDB--UNI\~AO remains close
to the proportionality diagonal, with \(R_C=1.050\), but its modest advantage
is enough to produce an inversion because its 49.22 percent vote share already
places it close to a majority. PP--PL is the conspicuous exception in Panel B:
with \(R_C=1.111\), it is the minimal connected coalition farthest from exact
proportionality. Its much larger representational advantage allows a coalition
with only 45.26 percent of the vote to reach 258 seats. The ideological
inversions thus include both modest departures that become consequential near
the majority threshold and a more pronounced case in which overrepresentation
compensates for a substantially larger vote deficit.


\section{Discussion}


The first result is a contrast between realized rarity and a broader range of
potential inversions. Among the 35 distinct cabinet party sets, only two invert.
Among the twenty minimal connected winning coalitions, seven do. Cabinet
participation therefore shows that coalition inversion was institutionally
realized, but only in a small number of the governing party configurations
actually observed. The ideological analysis shows, in turn, that these cabinets
do not exhaust the politically structured combinations under which the same
electoral distributions can reverse vote and seat majorities. The relaxation of
exact connectedness reinforces this interpretation. Allowing one represented
party inside an ideological span to be omitted preserves the strongest
inversions in 2014 and 2022 and reveals a stronger case in 2018, where a
coalition with 47.59 percent of the vote reaches exactly 257 seats\footnote{
Appendix~\ref{app:one-gap} compares the strongest inversions
under exact connectedness and the one-gap relaxation.}.


A second implication is that inversion is not simply a stronger form of
overrepresentation. Overrepresentation is common among cabinet coalitions: 33 of
the 35 sets have \(R_C>1\), yet only two combine a vote minority with a seat
majority. \(R_C\) measures the coalition's advantage relative to its vote share,
while \(d_C\) expresses that advantage in seats. Whether this advantage changes
majority status also depends on \(q_C\), which determines how far the coalition
lies from 257 seats under exact proportionality. A positive differential is
therefore insufficient for inversion unless it is large enough to bridge that
remaining gap to a seat majority. 



The Brazilian case adds a second stage to the aggregation problem identified
by Kurrild-Klitgaard and Miller. At the coalition level, positive and negative
party differentials reinforce or offset one another as parties are combined.
In Brazil, however, those party differentials are themselves national
aggregates of twenty-seven independently determined district outcomes. The
decomposition distinguishes the part of the coalition differential associated
with departures from proportionality within districts from the part arising
because the sum of district-level proportional quotas differs from the
coalition's national quota.  Both components carry into national party seat totals and, ultimately, into
coalition majority status.

Cabinet recomposition changes not only the parties represented in government but
also the distribution of their combined support across districts with different
electoral weights. The Dilma and Temer configurations draw on the same 2014
election, yet the former benefit from between-district weighting while the
latter are disadvantaged by it. Its contribution therefore depends on which
parties are combined rather than providing governing coalitions with a uniform
advantage.


The decomposition also places the inversion cases within a broader set of
coalition configurations. Some of the patterns observed
among inverted coalitions are not distinctive to inversion itself. All 35
distinct cabinet party sets have a positive within-district component,
whereas the between-district component is positive in only thirteen. The
two cabinet inversions therefore belong to a much wider pattern in which
cabinet composition repeatedly combines positive within-district
contributions, including in party sets that remain minorities in both votes
and seats or majorities in both. The party-level contributions help explain
this regularity: larger cabinet members often supply substantial positive
\(A_i\) contributions, while smaller parties can reinforce or offset them.


The ideological coalitions also show that similar majority outcomes can rest on
different representational advantages. MDB--UNIÃO benefits chiefly from the
territorial distribution of its support, whereas PP--PL obtains most of its
surplus through vote-to-seat translation within districts. The one-gap results
extend this contrast: favorable territorial weighting often provides the larger
contribution among the 2022 inversions, even though positive within-district
gains remain widespread. Component size must also be distinguished from its role
in reaching the majority threshold. A relatively small positive \(B_C\) can
provide the remaining seats needed after \(A_C\) has brought a coalition close
to 257, even when \(A_C\) accounts for most of the total differential. The
decomposition thus shows that coalitions with the same inversion status, and
even coalitions with similar overall representational advantages, can occupy
quite different positions with respect to within-district allocation and
between-district weighting.

This result deepens Miller's aggregation argument in an important respect. A
fixed electoral distribution does not imply a fixed political beneficiary of
coalition inversion. The same party vote and seat totals can sustain different
majority relations across governing configurations and can place potential
inversions in different regions of the party system. The Dilma inversion ends
when the remaining cabinet parties lose their seat majority, whereas the initial
Lula inversion ends when the expanded cabinet acquires a vote majority as well.
Neither transition reflects a new electoral allocation. The same party vote and
seat totals instead support different coalition-level majority relations at
different moments of the legislative term.

The ideological results show that this variation is not only temporal. In 2022,
the electoral allocation supports both Lula's broad, sparse, left-anchored
cabinet inversion and potential connected inversions on the right, including the
compact PP--PL interval\footnote{
  Figure~\ref{fig:cabinet-membership-inversions} in
  Appendix~\ref{app:cabinet-interval-bridge} displays the cabinet's membership
  in the election-year ideological order.}. These are politically different but
overlapping combinations: UNIÃO belongs to both the cabinet and PP--PL, carrying
its positive seat differential into each. The same party-level advantage can
therefore contribute to majorities with different political anchors. The 2018
election provides the converse pattern: no observed cabinet party set inverts,
even though an exact-connected inversion exists and the one-gap relaxation
reveals a stronger one. Party-level overrepresentation consequently does not
imply a stable advantage for one governing bloc or ideological side. Large
parties can contribute substantial positive or negative differentials, but which
political combinations benefit from them depends on how those contributions are
distributed across coalitions. The analysis identifies these representational
consequences of coalition structure; it does not imply that parties select
partners in order to exploit them.

Notice that cabinet membership identifies
the parties represented in government, not the coalition supporting the
executive in every legislative vote. Likewise, the vote share attributed to
a cabinet or potential coalition is obtained by summing the election-year
votes of its member parties; it is not a direct electoral endorsement of the
later coalition as a collective actor. The ideological analysis imposes a
different restriction: admissible coalitions are defined relative to a fixed
estimated election-year party order, so alternative party placements can
change the boundaries of connected or one-gap coalitions. Finally, the
\(A_C+B_C\) decomposition separates the observed coalition differential
into a within-district allocation component and a between-district weighting
component. It does not identify the independent causal effects of district
magnitude, thresholds, electoral pooling, malapportionment, turnout, or the
other institutional and geographical processes that contribute to those
components.





\section{Conclusion}

The Brazilian case shows that coalition inversion is neither merely a
theoretical possibility nor a routine feature of coalition politics under
proportional representation. It occurs among observed governing party sets, but
rarely: only two of the 35 distinct cabinet configurations invert. The same
electoral distributions, however, permit a broader range of inversions among
ideologically structured parliamentary majorities. Seven of the twenty minimal
connected winning coalitions invert, with cases in all three elections. The
one-gap analysis reveals further inversions beyond the exact-connected cases.
The central empirical result is therefore a contrast between realized rarity and
a wider set of politically structured inversion possibilities.

The institutional consequence is not simply that some parties are
overrepresented. Coalition composition determines which representational
advantages are combined into a majority and which are offset by partners'
disadvantages. The cabinet results show that governing combinations can
be advantaged or disadvantaged by between-district weighting
under the same electoral allocation.
The ideological results show that this allocation can make compact
majorities available in one part of the party order while requiring
broader combinations elsewhere. Coalition inversion identifies when these
differences in coalition composition become a reversal between vote and
seat majorities.

Proportional representation is usually described through the relationship
between each party's votes and seats. Once those allocations are combined
into parliamentary majorities, however, an additional structure becomes
visible. Coalition inversion identifies the configurations in which ordinary
party-level deviations are sufficient to reverse the relation between vote
majorities and seat majorities. This structure makes it possible to ask where
such reversals lie within a party system, whether they are realized in
governing coalitions, and how their incidence changes as the same legislature
is recombined. At the same time, the relation being reversed is a basic one
for democratic representation: a set of parties supported by less than half
of the national vote can control a majority of legislative seats. Coalition
inversions therefore reveal both an analytical property of proportional
systems and a majoritarian tension within the parliamentary outcomes those
systems can produce.



\section*{Statements and Declarations}  

\textbf{Funding.} No funding was received to assist with the preparation of this manuscript.

\textbf{Competing interests.} The author has no competing interests to declare that are relevant to the content of this article.

\textbf{Author contribution.} The author is the sole author and was responsible for conceptualization, methodology, formal analysis, software, investigation, data curation, visualization, and writing.

\textbf{Data availability.} For peer review, the data and replication materials are available in an anonymized repository: \url{https://anonymous.4open.science/r/electoral-inversion-B6B5/README.md}. For final publication, the analysis uses public election data from Brazil's Tribunal Superior Eleitoral and cabinet/party classification inputs documented in the replication repository.

\textbf{Code availability.} For peer review, the code necessary to reproduce the analysis is available in an anonymized repository: \url{https://anonymous.4open.science/r/electoral-inversion-B6B5/README.md}. For final publication, the replication code is available at the public repository after acceptance.

\textbf{Ethics approval.} This study uses publicly available aggregate electoral, cabinet, and party-classification data. No new human participants were recruited by the author. 

\textbf{Consent to participate.} Not applicable. This study uses publicly available aggregate and documentary data.

\textbf{Consent for publication.} Not applicable. This study uses publicly available aggregate and documentary data.

\textbf{AI assistance.} ChatGPT was used for copy-editing and to assist with extracting data from Wikipedia. All data used in the analysis were checked by the author.

% VERIFY before submission: all statements and declarations should be confirmed by the author.

\printbibliography


\clearpage
\appendix

\section{Supporting coalition data}
\label{app:accounting-diagnostics}

\subsection{Cabinet membership}
\label{app:composition}

% Generated by writing/make_appendix_assets.py; do not edit.
% source: processing/Processing/output/paper/raw/cabinet_party_sets.csv
\begingroup
\small
\setlength{\tabcolsep}{5pt}
\renewcommand{\arraystretch}{1.0}
\setlength{\LTpre}{4pt}
\setlength{\LTpost}{4pt}
\begin{longtable}{@{}L{0.08\textwidth}L{0.78\textwidth}r@{}}
\caption{Distinct election-year cabinet party sets}\label{tab:full-cabinet-composition}\\
\toprule
Set & Election-year member parties & Days \\
\midrule
\endfirsthead
\multicolumn{3}{@{}l}{\small Table \thetable\ (continued)}\\
\toprule
Set & Election-year member parties & Days \\
\midrule
\endhead
\midrule
\multicolumn{3}{r@{}}{\footnotesize Continued on next page}\\
\endfoot
\bottomrule
\endlastfoot
\addlinespace[4pt]\multicolumn{3}{@{}l}{\textit{Dilma Rousseff (2014 election)}}\\*
14-01 & PCdoB, PDT, PMDB, PP, PR, PRB, PROS, PSD, PT, PTB & 77 \\
14-02 & PCdoB, PDT, PMDB, PP, PR, PRB, PSD, PT, PTB & 365 \\
14-03 & PCdoB, PDT, PMDB, PP, PR, PROS, PSD, PT, PTB & 13 \\
14-04 & PCdoB, PDT, PMDB, PP, PR, PSD, PT, PTB & 14 \\
14-05 & PCdoB, PDT, PMDB, PR, PSD, PT, PTB & 5 \\
14-06 & PCdoB, PDT, PMDB, PR, PT, PTB & 23 \\
\addlinespace[4pt]\multicolumn{3}{@{}l}{\textit{Michel Temer (2014 election)}}\\*
14-07\(^{\dagger}\) & DEM, PMDB, PP, PPS, PR, PRB, PSB, PSD, PSDB, PTB, PV & 520 \\
14-08\(^{\dagger}\) & DEM, PMDB, PP, PPS, PR, PRB, PSD, PSDB, PTB, PV & 158 \\
14-09 & DEM, PMDB, PP, PR, PRB, PSD, PSDB, PTB, PV & 12 \\
14-10 & DEM, PMDB, PP, PRB, PSD, PSDB, PTB, PV & 4 \\
14-11 & PMDB, PP, PRB, PSD, PSDB, PTB, PV & 90 \\
14-12\(^{\dagger}\) & PMDB, PP, PRB, PSD, PSDB, PV & 180 \\
\addlinespace[4pt]\multicolumn{3}{@{}l}{\textit{Jair Bolsonaro (2018 election)}}\\*
18-01 & DEM, MDB, NOVO, PP, PSL & 406 \\
18-02\(^{\dagger}\) & DEM, MDB, NOVO, PP, PSDB, PSL & 8 \\
18-03\(^{\dagger}\) & DEM, NOVO, PP, PSDB, PSL & 78 \\
18-04 & DEM, PP, PSDB, PSL & 41 \\
18-05 & DEM, PP, PSD, PSDB, PSL & 1 \\
18-06 & DEM, PP, PSD, PSL & 93 \\
18-07 & DEM, PSD, PSL & 81 \\
18-08 & DEM, PSC, PSD, PSL & 65 \\
18-09\(^{\dagger}\) & DEM, PRB, PSC, PSD, PSL & 47 \\
18-10\(^{\dagger}\) & DEM, PR, PRB, PSC, PSD, PSL & 119 \\
18-11 & DEM, PP, PR, PRB, PSC, PSD, PSL & 239 \\
18-12 & DEM, PP, PR, PRB, PSC, PSL & 3 \\
18-13 & PP, PR, PSC & 1 \\
18-14 & PP, PR, PRB, PSC & 2 \\
18-15 & PP, PR, PRB & 1 \\
18-16 & PP, PRB, PSD & 41 \\
18-17\(^{\dagger}\) & DEM, PP, PRB, PSD, PSL & 233 \\
18-18\(^{\dagger}\) & DEM, NOVO, PRB, PSD, PSL & 2 \\
\addlinespace[4pt]\multicolumn{3}{@{}l}{\textit{Lula III (2022 election)}}\\*
22-01 & MDB, PCdoB, PDT, PSB, PSD, PSOL, PT, REDE, UNIÃO & 255 \\
22-02 & MDB, PCdoB, PDT, PP, PSB, PSD, PSOL, PT, REDE, REPUBLICANOS, UNIÃO & 807 \\
22-03 & MDB, PCdoB, PDT, PP, PSB, PSD, PSDB, PSOL, PT, REDE, REPUBLICANOS, UNIÃO & 5 \\
22-04 & MDB, PCdoB, PDT, PP, PSB, PSD, PSDB, PSOL, PT, REDE, REPUBLICANOS & 9 \\
22-05 & MDB, PCdoB, PDT, PP, PSB, PSD, PSOL, PT, REDE, REPUBLICANOS & 98 \\
\end{longtable}
\noindent\footnotesize\textit{Note:} Days sum actual covered dates across appearances, not first-to-last spans. \(\dagger\): set includes provisionally coded days. Full histories and coding details remain in the replication files.
\endgroup


\subsection{Baseline minimal winning intervals}
\label{app:minimal-connected-winning}

% Generated by writing/make_appendix_assets.py; do not edit.
% source: processing/Processing/output/paper/tables/table_appendix_minimal_connected_winning_intervals.csv
\begin{table}[!htbp]
\centering\small
\caption{All baseline minimal connected winning intervals}
\label{tab:minimal-connected-winning-intervals}
\setlength{\tabcolsep}{8pt}
\renewcommand{\arraystretch}{1.0}
\begin{tabular}{@{}llrrr@{}}
\toprule
Election & Interval endpoints & Vote \% & Seats & \(d_C\) \\
\midrule
2014 & PCdoB--PMDB & 50.17 & 265 & 7.63 \\
2014 & PT--PSD & 54.34 & 291 & 12.26 \\
2014 & PDT--PSDB & 51.80 & 276 & 10.25 \\
2014 & PSB--PTN\(^{*}\) & 48.92 & 260 & 9.03 \\
2014 & PTB--PR\(^{*}\) & 47.45 & 257 & 13.60 \\
2014 & PT DO B--PSDC\(^{*}\) & 48.84 & 257 & 6.44 \\
2014 & SOLIDARIEDADE--PSL\(^{*}\) & 48.84 & 257 & 6.46 \\
2014 & PMDB--PP & 51.23 & 272 & 9.19 \\
\addlinespace[4pt]
2018 & PT--PSDB\(^{*}\) & 49.95 & 267 & 10.75 \\
2018 & PDT--PR & 50.59 & 268 & 8.47 \\
2018 & PSB--PRB & 51.02 & 269 & 7.26 \\
2018 & SOLIDARIEDADE--PSL & 50.12 & 262 & 4.86 \\
2018 & PMN--NOVO & 50.93 & 257 & -4.29 \\
2018 & MDB--PP & 54.34 & 285 & 6.24 \\
2018 & PSD--PSC & 50.55 & 258 & -1.34 \\
2018 & PSDB--DEM & 50.82 & 257 & -3.72 \\
\addlinespace[4pt]
2022 & PT--PP & 57.18 & 284 & -9.36 \\
2022 & PSB--PSC & 53.82 & 261 & -15.10 \\
2022 & MDB--UNIÃO\(^{*}\) & 49.22 & 265 & 12.51 \\
2022 & PP--PL\(^{*}\) & 45.26 & 258 & 25.81 \\
\bottomrule
\end{tabular}
\par\vspace{4pt}\begin{minipage}{\textwidth}\footnotesize\textit{Note:} Each interval contains all seat-winning parties between its endpoints in the Figure~\ref{fig:minimal-connected-inversions} election-year order. \(*\): inversion, classified before rounding; vote shares use all valid federal-deputy votes.
\end{minipage}
\end{table}


\subsection{Selected member-party contributions}
\label{app:party-contributions}

Here \(d_i=A_i+B_i\), \(A_C=\sum_{i\in C}A_i\), and
\(B_C=\sum_{i\in C}B_i\). These ex post accounting contributions do not
identify separate causal effects; in 2014 and 2018, party attribution also
reflects joint-list seat allocation.

% Generated by writing/make_appendix_assets.py; do not edit.
% source: processing/Processing/output/paper/tables/table_coalition_party_component_extremes.csv
\begin{table}[!htbp]
\centering\small
\caption{Selected member-party contributions}
\label{tab:coalition-party-component-extremes}
\setlength{\tabcolsep}{4pt}
\renewcommand{\arraystretch}{1.15}
\begin{tabularx}{\textwidth}{@{}L{0.20\textwidth}*{4}{>{\raggedright\arraybackslash}X}@{}}
\toprule
Case & \multicolumn{2}{c}{Within-district \(A_i\)} & \multicolumn{2}{c}{Between-district \(B_i\)} \\
\cmidrule(lr){2-3}\cmidrule(l){4-5}
 & Positive & Negative & Positive & Negative \\
\midrule
Cabinet\newline 2014/14-05 & \mbox{PR} (5.04)\newline \mbox{PSD} (4.65) & \mbox{PT} (-1.22)\newline \mbox{PCdoB} (-0.57) & \mbox{PMDB} (4.22)\newline \mbox{PDT} (1.56) & \mbox{PT} (-1.21)\newline \mbox{PR} (-0.74)\\ \addlinespace[3pt]
Cabinet\newline 2022/22-01 & \mbox{UNIÃO} (9.36)\newline \mbox{PT} (9.07) & \mbox{PSB} (-5.55)\newline \mbox{PSOL} (-3.40) & \mbox{UNIÃO} (1.75)\newline \mbox{PDT} (1.02) & \mbox{PSOL} (-2.66)\newline \mbox{PT} (-2.13)\\ \addlinespace[3pt]
Baseline\newline 2022/MDB--UNIÃO & \mbox{UNIÃO} (9.36)\newline \mbox{MDB} (5.12) & \mbox{PTB} (-5.94)\newline \mbox{PODE} (-5.66) & \mbox{PP} (3.17)\newline \mbox{REPUBLICANOS} (2.94) & \mbox{PSDB} (-0.54)\newline \mbox{MDB} (-0.40)\\ \addlinespace[3pt]
\bottomrule
\end{tabularx}
\par\vspace{4pt}\begin{minipage}{\textwidth}\footnotesize\textit{Note:} Entries are the two largest positive and two most negative member contributions per component, in seats. Baseline denotes the seat-winning \(k=0\) interval. PODE is the election-year label for Podemos.
\end{minipage}
\end{table}


\clearpage
\subsection{Inverted cabinets in ideological order}
\label{app:cabinet-interval-bridge}

Figure~\ref{fig:cabinet-membership-inversions} shows the two cabinet inversions
as selective memberships in the same ordinal party orders used in
Figure~\ref{fig:minimal-connected-inversions}. Both include PT and span a wide
part of their election-year order, omitting twelve represented parties inside
the 14-05 span and eleven inside the 22-01 span. In particular, Lula's
PSOL--UNIÃO span is broad, sparse, and left-anchored, whereas the 2022 minimal
connected inversions lie entirely to PT's right. This comparison concerns the
two displayed cabinets.

% PROVENANCE-BEGIN appendix-selected-cabinet-spans
% note: Compact operands: build/results/manuscript_results.csv; block=appendix-selected-cabinet-spans
% source: build/results/domains/tables/table_appendix_cabinet_interval_bridge.csv
% row:
%   key: display_label=14-05
%   row_at_generation: 5
%   fields:
%     cabinet_parties: PCdoB, PDT, PMDB, PR, PSD, PT, PTB
%     closure_start_party: PCdoB
%     closure_end_party: PR
%     closure_gap_n: 12
% source: build/results/domains/tables/table_appendix_cabinet_interval_bridge.csv
% row:
%   key: display_label=22-01
%   row_at_generation: 31
%   fields:
%     cabinet_parties: MDB, PCdoB, PDT, PSB, PSD, PSOL, PT, REDE, UNIÃO
%     closure_start_party: PSOL
%     closure_end_party: UNIÃO
%     closure_gap_n: 11
% PROVENANCE-END appendix-selected-cabinet-spans

\begin{figure}[!htbp]
\centering
\includegraphics[width=\textwidth]{../../../build/assets/cabinet_membership_inversions.pdf}
\caption{Cabinet inversions in election-year ideological order. Filled markers
identify members; open circles identify nonmembers, including omitted parties
inside each cabinet's span. PT is a filled diamond. All 28 seat-winning parties
in 2014 and all 23 in 2022 appear at equally spaced ordinal positions. Dates
are inclusive observed intervals; votes and seats use the corresponding election.}
\label{fig:cabinet-membership-inversions}
\end{figure}

\clearpage
\section{Robustness of the ideological coalition domain}
\label{app:ideological-robustness}

\subsection{Admissible domains and the one-gap relaxation}
\label{app:one-gap}

Let \(p_1,\ldots,p_n\) order the parties in \(P^*\), where \(P^*=P^+\)
contains seat-winning parties in the primary specification and \(P^*=P\)
contains all parties in the sensitivity analysis. For nonempty
\(C\subseteq P^*\), let \(a=\min\{j:p_j\in C\}\) and
\(b=\max\{j:p_j\in C\}\). Define
\[
\operatorname{span}(C)=\{p_a,\ldots,p_b\},\qquad
 g(C)=|\operatorname{span}(C)\setminus C|,\qquad
 \mathcal D_k=\{\varnothing\ne C\subseteq P^*:g(C)\leq k\}.
\]
Thus, \(k=0\) requires connectedness, while \(k=1\) permits at most one
omitted party inside the span. A winning \(C\) is minimal only if no proper
subset admissible in the same \(\mathcal D_k\) retains at least 257 seats.
Minimality is recomputed for each \(k\) and party universe: nested admissible
domains do not imply nested collections of minimal winning coalitions.

All twenty baseline minimal majorities remain admissible and winning under
\(k=1\), but only six remain minimal. Table~\ref{tab:one-gap-summary}
compares the counts. The  one-gap domain contains 100 inversions among
234 minimal winning coalitions. Full enumeration, memberships, omissions, and
decompositions for both universes remain in the machine-readable replication
outputs.

% PROVENANCE-BEGIN appendix-compact-domain-totals
% note: Compact operands: build/results/manuscript_results.csv; block=appendix-compact-domain-totals
% source: build/results/accounting/tables/prose_analysis_summaries.csv
% row:
%   key: summary=seat-winning-k0-totals; metric=minimal_seat_majority_coalitions; aggregation=sum
%   row_at_generation: 3
%   fields:
%     value: 20
% source: build/results/accounting/tables/prose_analysis_summaries.csv
% row:
%   key: summary=seat-winning-k1-surviving-connected; metric=count; aggregation=sum
%   row_at_generation: 5
%   fields:
%     value: 6
% source: build/results/domains/tables/ideological_universe_comparison.csv
% row:
%   key: election=2014; ideological_universe=seat_winning; k=1
%   row_at_generation: 3
%   fields:
%     minimal_seat_majority_coalitions: 87
%     minimal_inversions: 46
% source: build/results/domains/tables/ideological_universe_comparison.csv
% row:
%   key: election=2018; ideological_universe=seat_winning; k=1
%   row_at_generation: 7
%   fields:
%     minimal_seat_majority_coalitions: 108
%     minimal_inversions: 42
% source: build/results/domains/tables/ideological_universe_comparison.csv
% row:
%   key: election=2022; ideological_universe=seat_winning; k=1
%   row_at_generation: 11
%   fields:
%     minimal_seat_majority_coalitions: 39
%     minimal_inversions: 12
% PROVENANCE-END appendix-compact-domain-totals

% Generated by writing/make_appendix_assets.py; do not edit.
% source: processing/Processing/output/paper/tables/ideological_universe_comparison.csv
\begin{table}[!htbp]
\centering\small
\caption{Minimal winning coalitions by ideological domain and party universe}
\label{tab:one-gap-summary}
\label{tab:ideological-universe-comparison}
\setlength{\tabcolsep}{6pt}
\renewcommand{\arraystretch}{1.12}
\begin{tabular}{@{}lcrrrr@{}}
\toprule
 & & \multicolumn{2}{c}{Seat-winning parties} & \multicolumn{2}{c}{All parties} \\
\cmidrule(lr){3-4}\cmidrule(l){5-6}
Election & \(k\) & Minimal majorities & Inversions & Minimal majorities & Inversions \\
\midrule
2014 & 0 & 8 & 4 & 8 & 4\\
2014 & 1 & 87 & 46 & 88 & 43\\
\addlinespace[3pt]
2018 & 0 & 8 & 1 & 8 & 0\\
2018 & 1 & 108 & 42 & 118 & 24\\
\addlinespace[3pt]
2022 & 0 & 4 & 2 & 4 & 2\\
2022 & 1 & 39 & 12 & 53 & 17\\
\bottomrule
\end{tabular}
\par\vspace{4pt}\begin{minipage}{\textwidth}\footnotesize\textit{Note:} Minimality is recomputed for each \(k\) and universe. The primary \(k=1\) rows sum to 234 minimal majorities and 100 inversions. The \(k=0\) and \(k=1\) rows are overlapping domains and are not pooled. Every vote share uses all valid federal-deputy votes in \(V\).
\end{minipage}
\end{table}


The strongest inversion is the one with the lowest national vote share among
the relevant minimal inversions. Under the  one-gap relaxation, the
strongest 2014 and 2022 cases retain their baseline member sets and vote--seat
outcomes: PTB--PR has 47.45 percent and 257 seats; PP--PL has 45.26 percent and
258 seats. The strongest 2018 case is PCdoB--PODE omitting PSDB, with
47.59 percent of all valid votes and 257 seats.

% PROVENANCE-BEGIN appendix-compact-one-gap-strongest
% note: Compact operands: build/results/manuscript_results.csv; block=appendix-compact-one-gap-strongest
% source: build/results/domains/tables/ideological_universe_comparison.csv
% row:
%   key: election=2014; ideological_universe=seat_winning; k=1
%   row_at_generation: 3
%   fields:
%     strongest_inversion_coalition: PTB--PR
%     strongest_inversion_vote_share: 0.4744616510767348
%     strongest_inversion_seats: 257
% source: build/results/domains/tables/ideological_universe_comparison.csv
% row:
%   key: election=2018; ideological_universe=seat_winning; k=1
%   row_at_generation: 7
%   fields:
%     strongest_inversion_coalition: PCdoB--PODE, omitting PSDB
%     strongest_inversion_vote_share: 0.4759116927539931
%     strongest_inversion_seats: 257
% source: build/results/domains/tables/ideological_universe_comparison.csv
% row:
%   key: election=2022; ideological_universe=seat_winning; k=1
%   row_at_generation: 11
%   fields:
%     strongest_inversion_coalition: PP--PL
%     strongest_inversion_vote_share: 0.4526056142903306
%     strongest_inversion_seats: 258
% PROVENANCE-END appendix-compact-one-gap-strongest

Among the 100  one-gap minimal inversions, \(A_C>0\) in 97 cases and
\(A_C>B_C\) in 77: 40 of 46 in 2014, 32 of 42 in 2018, and 5 of 12 in 2022.
The three negative-\(A_C\) cases (two in 2014 and one in 2022) have positive
between-district components that more than offset the shortfall.

% PROVENANCE-BEGIN appendix-compact-component-frequencies
% note: Compact operands: build/results/manuscript_results.csv; block=appendix-compact-component-frequencies
% note: 97 positive cases = 100 minimal inversions minus 3 negative cases; absence of zeros and exact signs are checked by writing/make_appendix_assets.py.
% source: build/results/accounting/tables/prose_analysis_summaries.csv
% row:
%   key: summary=ideology-component-dominance; metric=minimal_inversions; aggregation=count; ideological_universe=seat_winning; k=1; election=all
%   row_at_generation: 61
%   fields:
%     value: 100
% source: build/results/accounting/tables/prose_analysis_summaries.csv
% row:
%   key: summary=ideology-component-dominance; metric=A_C_gt_B_C; aggregation=count; ideological_universe=seat_winning; k=1; election=all
%   row_at_generation: 62
%   fields:
%     value: 77
% source: build/results/accounting/tables/prose_analysis_summaries.csv
% row:
%   key: summary=ideology-component-dominance; metric=minimal_inversions; aggregation=count; ideological_universe=seat_winning; k=1; election=2014
%   row_at_generation: 68
%   fields:
%     value: 46
% source: build/results/accounting/tables/prose_analysis_summaries.csv
% row:
%   key: summary=ideology-component-dominance; metric=A_C_gt_B_C; aggregation=count; ideological_universe=seat_winning; k=1; election=2014
%   row_at_generation: 69
%   fields:
%     value: 40
% source: build/results/accounting/tables/prose_analysis_summaries.csv
% row:
%   key: summary=ideology-component-dominance; metric=minimal_inversions; aggregation=count; ideological_universe=seat_winning; k=1; election=2018
%   row_at_generation: 75
%   fields:
%     value: 42
% source: build/results/accounting/tables/prose_analysis_summaries.csv
% row:
%   key: summary=ideology-component-dominance; metric=A_C_gt_B_C; aggregation=count; ideological_universe=seat_winning; k=1; election=2018
%   row_at_generation: 76
%   fields:
%     value: 32
% source: build/results/accounting/tables/prose_analysis_summaries.csv
% row:
%   key: summary=ideology-component-dominance; metric=minimal_inversions; aggregation=count; ideological_universe=seat_winning; k=1; election=2022
%   row_at_generation: 82
%   fields:
%     value: 12
% source: build/results/accounting/tables/prose_analysis_summaries.csv
% row:
%   key: summary=ideology-component-dominance; metric=A_C_gt_B_C; aggregation=count; ideological_universe=seat_winning; k=1; election=2022
%   row_at_generation: 83
%   fields:
%     value: 5
% source: build/results/accounting/tables/prose_analysis_summaries.csv
% row:
%   key: summary=seat-winning-k1-negative-within; metric=A_C; aggregation=count_negative
%   row_at_generation: 7
%   fields:
%     value: 3
% source: build/results/accounting/tables/prose_analysis_summaries.csv
% row:
%   key: summary=seat-winning-2014-k1-negative-within; metric=A_C; aggregation=count_negative
%   row_at_generation: 8
%   fields:
%     value: 2
% source: build/results/accounting/tables/prose_analysis_summaries.csv
% row:
%   key: summary=seat-winning-2022-k1-negative-within; metric=A_C; aggregation=count_negative
%   row_at_generation: 9
%   fields:
%     value: 1
% PROVENANCE-END appendix-compact-component-frequencies

\newpage
\subsection{Sensitivity to the ideological party universe}
\label{app:ideological-universes}

Restoring parties without seats removes the 2018 exact-connected PT--PSDB
inversion: its interval includes PMB and acquires a vote majority. The strongest
2014 and 2022 endpoint regions persist. Under all-party \(k=0\), the strongest
cases remain PTB--PR in 2014, with 47.59 percent of the vote and 257 seats, and
PP--PL in 2022, with 45.35 percent and 258 seats. Among the six all-party
\(k=0\) inversions, \(A_C>0\) in five; MDB--UNIÃO is the sole exception.
Under all-party \(k=1\), their strongest
cases omit PPL and DC, respectively, recovering the  PTB--PR and PP--PL
memberships and outcomes. The strongest 2018 all-party one-gap case is
PCdoB--PODE omitting PSDB, with 47.82 percent of the vote and the same 257 seats,
compared with 47.59 percent in the seat-winning one-gap specification.

% PROVENANCE-BEGIN appendix-compact-universe-strongest
% note: Compact operands: build/results/manuscript_results.csv; block=appendix-compact-universe-strongest
% source: build/results/domains/tables/ideological_universe_comparison.csv
% row:
%   key: election=2014; ideological_universe=all_parties; k=0
%   row_at_generation: 2
%   fields:
%     strongest_inversion_coalition: PTB--PR
%     strongest_inversion_vote_share: 0.47591256254894826
%     strongest_inversion_seats: 257
% source: build/results/domains/tables/ideological_universe_comparison.csv
% row:
%   key: election=2022; ideological_universe=all_parties; k=0
%   row_at_generation: 10
%   fields:
%     strongest_inversion_coalition: PP--PL
%     strongest_inversion_vote_share: 0.4534989317772354
%     strongest_inversion_seats: 258
% source: build/results/accounting/tables/prose_analysis_summaries.csv
% row:
%   key: summary=all-parties-k0-totals; metric=minimal_inversions; aggregation=sum; ideological_universe=all_parties; k=0; election=all
%   row_at_generation: 4
%   fields:
%     value: 6
%     source_row_count: 3
% source: build/results/accounting/tables/prose_analysis_summaries.csv
% row:
%   key: summary=all-parties-k0-positive-within; metric=A_C; aggregation=count_positive; ideological_universe=all_parties; k=0; election=all
%   row_at_generation: 10
%   fields:
%     value: 5
%     source_row_count: 6
% source: build/results/domains/raw/ideology_k_gap_minimal_accounting_both_universes.csv
% row:
%   key: election=2014; ideological_universe=all_parties; k=0; left_endpoint=PSB; right_endpoint=PTN
%   row_at_generation: 258
%   fields:
%     votes: 47627762
%     national_vote_total: 97355354
%     seats: 260
%     A_C: 10.508783851460613
%     inversion: true
% source: build/results/domains/raw/ideology_k_gap_minimal_accounting_both_universes.csv
% row:
%   key: election=2014; ideological_universe=all_parties; k=0; left_endpoint=PTB; right_endpoint=PR
%   row_at_generation: 259
%   fields:
%     votes: 46332636
%     national_vote_total: 97355354
%     seats: 257
%     A_C: 11.620538081371619
%     inversion: true
% source: build/results/domains/raw/ideology_k_gap_minimal_accounting_both_universes.csv
% row:
%   key: election=2014; ideological_universe=all_parties; k=0; left_endpoint=PT DO B; right_endpoint=PSDC
%   row_at_generation: 260
%   fields:
%     votes: 47691320
%     national_vote_total: 97355354
%     seats: 257
%     A_C: 6.2665828795721277
%     inversion: true
% source: build/results/domains/raw/ideology_k_gap_minimal_accounting_both_universes.csv
% row:
%   key: election=2014; ideological_universe=all_parties; k=0; left_endpoint=SOLIDARIEDADE; right_endpoint=PSL
%   row_at_generation: 261
%   fields:
%     votes: 47687533
%     national_vote_total: 97355354
%     seats: 257
%     A_C: 6.5862104528808336
%     inversion: true
% source: build/results/domains/raw/ideology_k_gap_minimal_accounting_both_universes.csv
% row:
%   key: election=2022; ideological_universe=all_parties; k=0; left_endpoint=MDB; right_endpoint=UNIÃO
%   row_at_generation: 479
%   fields:
%     votes: 54687303
%     national_vote_total: 109413508
%     seats: 265
%     A_C: -1.2923532538854881
%     inversion: true
% source: build/results/domains/raw/ideology_k_gap_minimal_accounting_both_universes.csv
% row:
%   key: election=2022; ideological_universe=all_parties; k=0; left_endpoint=PP; right_endpoint=PL
%   row_at_generation: 480
%   fields:
%     votes: 49618909
%     national_vote_total: 109413508
%     seats: 258
%     A_C: 22.820505594331628
%     inversion: true
% source: build/results/domains/tables/ideological_universe_comparison.csv
% row:
%   key: election=2014; ideological_universe=all_parties; k=1
%   row_at_generation: 4
%   fields:
%     strongest_inversion_coalition: PTB--PR, omitting PPL
%     strongest_inversion_vote_share: 0.4744616510767348
%     strongest_inversion_seats: 257
% source: build/results/domains/tables/ideological_universe_comparison.csv
% row:
%   key: election=2018; ideological_universe=all_parties; k=1
%   row_at_generation: 8
%   fields:
%     strongest_inversion_coalition: PCdoB--PODE, omitting PSDB
%     strongest_inversion_vote_share: 0.47823504167693237
%     strongest_inversion_seats: 257
% source: build/results/domains/tables/ideological_universe_comparison.csv
% row:
%   key: election=2022; ideological_universe=all_parties; k=1
%   row_at_generation: 12
%   fields:
%     strongest_inversion_coalition: PP--PL, omitting DC
%     strongest_inversion_vote_share: 0.4526056142903306
%     strongest_inversion_seats: 258
% source: build/results/domains/raw/ideology_k_gap_minimal_accounting_both_universes.csv
% row:
%   key: election=2018; ideological_universe=all_parties; k=0; left_endpoint=PT; right_endpoint=PSDB
%   row_at_generation: 351
%   fields:
%     parties: PT, PDT, PSB, REDE, PPS, PV, PTB, AVANTE, SOLIDARIEDADE, PMN, PMB, PHS, MDB, PSD, PSDB
%     vote_share: 0.5018337809531631
%     inversion: false
% PROVENANCE-END appendix-compact-universe-strongest

The national denominator \(V\) remains all valid federal-deputy votes in every
specification. Changing the ideological universe changes membership, adjacency,
coalition vote numerators, and decomposition totals; identical endpoints need
not identify identical member sets. In particular, baseline MDB--UNIÃO has
\(A_C=2.95\) and \(B_C=9.55\), whereas the all-party interval includes additional
seatless members and has \(A_C=-1.29\) and \(B_C=9.88\). It remains inverted,
with 49.98 percent of all valid votes and 265 seats.

% PROVENANCE-BEGIN appendix-compact-mdb-uniao
% note: Compact operands: build/results/manuscript_results.csv; block=appendix-compact-mdb-uniao
% source: build/results/domains/raw/ideology_k_gap_minimal_accounting_both_universes.csv
% row:
%   key: election=2022; ideological_universe=seat_winning; k=0; left_endpoint=MDB; right_endpoint=UNIÃO
%   row_at_generation: 214
%   fields:
%     parties: MDB, PSDB, PSD, PODE, PROS, PTB, PP, REPUBLICANOS, PSC, UNIÃO
%     vote_share: 0.492189812614362
%     seats: 265
%     A_C: 2.9531634698279121
%     B_C: 9.5534626590043707
%     inversion: true
% source: build/results/domains/raw/ideology_k_gap_minimal_accounting_both_universes.csv
% row:
%   key: election=2022; ideological_universe=all_parties; k=0; left_endpoint=MDB; right_endpoint=UNIÃO
%   row_at_generation: 479
%   fields:
%     parties: MDB, PMN, PSDB, PSD, PMB, PODE, PROS, PRTB, AGIR, PTB, PP, DC, REPUBLICANOS, PSC, UNIÃO
%     vote_share: 0.49982222487556105
%     seats: 265
%     A_C: -1.2923532538854881
%     B_C: 9.8835518927226591
%     inversion: true
% PROVENANCE-END appendix-compact-mdb-uniao

\end{document}
